\section{团队详情（41人）}
\subsection{领导团队（16人）}
\vspace{4pt}
\noindent{\bfseries Sabrina Maniscalco}\quad{\color{graytext}\footnotesize CEO and Co-founder}
\par{\footnotesize\color{graytext}PhD in Physics - Quantum Information Physics (2004) | University of Palermo, University of Helsinki \textbar\  h-index: 53 \textbar\  \href{https://www.linkedin.com/in/sabrina-maniscalco-0402642/}{LinkedIn} \textbar\  \href{https://scholar.google.com/citations?user=EZ0GPwgAAAAJ}{Scholar}}
\par{\footnotesize\color{mainred}\textbf{经历}}
\begin{itemize}[leftmargin=1.2em,itemsep=0pt,topsep=0pt]
\item {\footnotesize University of Helsinki: Professor of Quantum Information, Computing and Logic (2020-present)}
\item {\footnotesize University of Turku: Chair of Theoretical Physics (2014-2020)}
\item {\footnotesize Heriot-Watt University: Associate Professor (Prior to 2014)}
\item {\footnotesize University of Sofia: Postdoctoral Researcher (2004-2006)}
\end{itemize}
\par{\footnotesize\color{mainred}\textbf{研究方向}}
\begin{itemize}[leftmargin=1.2em,itemsep=0pt,topsep=0pt]
\item {\footnotesize 开放量子系统与非马尔可夫动力学}
\item {\footnotesize 量子退相干与纠缠}
\item {\footnotesize 近期量子设备上的量子模拟}
\item {\footnotesize 量子错误缓解}
\item {\footnotesize 量子信息理论}
\end{itemize}
\par{\footnotesize\color{mainred}\textbf{代表论文}}
\begin{itemize}[leftmargin=1.2em,itemsep=0pt,topsep=0pt]
\item {\footnotesize Sudden transition between classical and quantum decoherence (2010, 641 citations)}
\item {\footnotesize Protecting entanglement via the quantum Zeno effect (2008, 550 citations)}
\end{itemize}
\par{\footnotesize\color{mainred}\textbf{成就}}
\begin{itemize}[leftmargin=1.2em,itemsep=0pt,topsep=0pt]
\item {\footnotesize 世界经济论坛技术先锋（2024年，与Algorithmiq共同获得）}
\item {\footnotesize Wellcome Leap Q4Bio 200万美元奖金得主（2026年）——唯一提供用于药物发现的可扩展端到端量子框架的团队}
\item {\footnotesize 芬兰量子技术卓越中心副主任}
\end{itemize}
\par{\footnotesize\color{mainred}\textbf{角色}} {\footnotesize CEO兼联合创始人。领导公司战略、融资和合作伙伴关系（IBM、Cleveland Clinic、CERN、NVIDIA）。推动将量子算法应用于生命科学和药物发现的使命。}
\vspace{2pt}\hrule\vspace{2pt}
\vspace{4pt}
\noindent{\bfseries Guillermo García-Pérez}\quad{\color{graytext}\footnotesize CSO and Co-founder}
\par{\footnotesize\color{graytext}PhD in Physics - Complex Systems and Network Science (2018) | University of Barcelona, University of Turku \textbar\  h-index: 18-25 \textbar\  \href{https://www.linkedin.com/in/guillermo-garc%C3%ADa-p%C3%A9rez-157143b9/}{LinkedIn}}
\par{\footnotesize\color{mainred}\textbf{经历}}
\begin{itemize}[leftmargin=1.2em,itemsep=0pt,topsep=0pt]
\item {\footnotesize University of Turku: Senior Researcher, Turku Centre for Quantum Physics (2018-2020)}
\item {\footnotesize University of Helsinki: Academy of Finland Grant Researcher (2021)}
\item {\footnotesize University of Barcelona: PhD Researcher, Dept. de Física Fonamental (2014-2018)}
\end{itemize}
\par{\footnotesize\color{mainred}\textbf{研究方向}}
\begin{itemize}[leftmargin=1.2em,itemsep=0pt,topsep=0pt]
\item {\footnotesize 复杂网络与几何重整化}
\item {\footnotesize 量子信息科学}
\item {\footnotesize 用于错误缓解的张量网络方法}
\item {\footnotesize 信息完备测量（IC-POVMs）}
\item {\footnotesize 用于生命科学与药物发现的量子算法}
\end{itemize}
\par{\footnotesize\color{mainred}\textbf{代表论文}}
\begin{itemize}[leftmargin=1.2em,itemsep=0pt,topsep=0pt]
\item {\footnotesize Multiscale unfolding of real networks by geometric renormalization (Nature Physics, 2018)}
\item {\footnotesize Learning to Measure: Adaptive Informationally Complete Generalized Measurements for Quantum Algorithms (PRX Quantum, 2021, 87 citations)}
\end{itemize}
\par{\footnotesize\color{mainred}\textbf{成就}}
\begin{itemize}[leftmargin=1.2em,itemsep=0pt,topsep=0pt]
\item {\footnotesize 共同开发了 Algorithmiq 的 TEM（张量网络错误缓解）算法}
\item {\footnotesize Algorithmiq Aurora 平台的核心架构师}
\item {\footnotesize 在《Nature Physics》和《Nature Communications》上发表了文章}
\end{itemize}
\par{\footnotesize\color{mainred}\textbf{角色}} {\footnotesize 首席科学官兼联合创始人。领导科学战略，开发量子算法（TEM、Aurora 平台），监督管理与 IBM Quantum 及其他硬件提供商的研究合作。}
\vspace{2pt}\hrule\vspace{2pt}
\vspace{4pt}
\noindent{\bfseries Matteo Rossi}\quad{\color{graytext}\footnotesize CTO and Co-founder}
\par{\footnotesize\color{graytext}PhD in Physics - Quantum Information and Quantum Simulation (2017) | University of Milan, University of Parma \textbar\  h-index: 25 \textbar\  \href{https://www.linkedin.com/in/matteo-rossi-9ab6b3183/}{LinkedIn}}
\par{\footnotesize\color{mainred}\textbf{经历}}
\begin{itemize}[leftmargin=1.2em,itemsep=0pt,topsep=0pt]
\item {\footnotesize University of Helsinki: Docent (Adjunct Professor) in Theoretical Physics (2020-present)}
\item {\footnotesize University of Turku: Researcher (2017-2020)}
\item {\footnotesize University of Nottingham: Researcher (2015-2017)}
\end{itemize}
\par{\footnotesize\color{mainred}\textbf{研究方向}}
\begin{itemize}[leftmargin=1.2em,itemsep=0pt,topsep=0pt]
\item {\footnotesize 量子信息理论与应用}
\item {\footnotesize 近期设备上的量子模拟}
\item {\footnotesize 错误缓解技术}
\item {\footnotesize 开放量子系统与量子计量学}
\item {\footnotesize 复杂网络上的量子行走}
\end{itemize}
\par{\footnotesize\color{mainred}\textbf{代表论文}}
\begin{itemize}[leftmargin=1.2em,itemsep=0pt,topsep=0pt]
\item {\footnotesize Learning to Measure: Adaptive Informationally Complete Generalized Measurements for Quantum Algorithms (PRX Quantum, 2021)}
\item {\footnotesize IBM Q Experience as a versatile experimental testbed for simulating open quantum systems}
\end{itemize}
\par{\footnotesize\color{mainred}\textbf{成就}}
\begin{itemize}[leftmargin=1.2em,itemsep=0pt,topsep=0pt]
\item {\footnotesize 共同开发了Algorithmiq关于信息完备测量的核心知识产权}
\item {\footnotesize 担任Algorithmiq技术平台的软件负责人}
\item {\footnotesize 在111篇出版物中被引用超过2,630次}
\end{itemize}
\par{\footnotesize\color{mainred}\textbf{角色}} {\footnotesize 首席技术官兼联合创始人。领导Algorithmiq量子软件平台的技术开发，监督工程团队和产品架构。}
\vspace{2pt}\hrule\vspace{2pt}
\vspace{4pt}
\noindent{\bfseries Boris Sokolov}\quad{\color{graytext}\footnotesize Lead Researcher and Co-founder}
\par{\footnotesize\color{graytext}PhD in Theoretical and Mathematical Physics - Open Quantum Systems and Quantum Networks (2020) | University of Turku, University of Helsinki \textbar\  h-index: 8-12}
\par{\footnotesize\color{mainred}\textbf{经历}}
\begin{itemize}[leftmargin=1.2em,itemsep=0pt,topsep=0pt]
\item {\footnotesize University of Helsinki: Researcher, QTF Centre of Excellence (2020-present)}
\item {\footnotesize University of Turku: PhD Researcher, Complex Systems Research Group (2016-2020)}
\end{itemize}
\par{\footnotesize\color{mainred}\textbf{研究方向}}
\begin{itemize}[leftmargin=1.2em,itemsep=0pt,topsep=0pt]
\item {\footnotesize 开放量子系统与相对论量子信息}
\item {\footnotesize 复杂量子网络与多体物理}
\item {\footnotesize 自适应信息完备测量}
\item {\footnotesize 变分量子算法}
\item {\footnotesize 量子自旋链中的纠缠结构}
\end{itemize}
\par{\footnotesize\color{mainred}\textbf{代表论文}}
\begin{itemize}[leftmargin=1.2em,itemsep=0pt,topsep=0pt]
\item {\footnotesize Learning to Measure: Adaptive Informationally Complete Generalized Measurements for Quantum Algorithms (PRX Quantum, 2021, 87 citations)}
\item {\footnotesize Emergent entanglement structures and self-similarity in quantum spin chains (Phil. Trans. R. Soc. A, 2022)}
\end{itemize}
\par{\footnotesize\color{mainred}\textbf{成就}}
\begin{itemize}[leftmargin=1.2em,itemsep=0pt,topsep=0pt]
\item {\footnotesize Algorithmiq 联合创始人（2020年）}
\item {\footnotesize Algorithmiq 测量优化和错误缓解知识产权的主要贡献者}
\item {\footnotesize 博士论文探讨了开放量子系统、相对论与复杂网络之间的联系}
\end{itemize}
\par{\footnotesize\color{mainred}\textbf{角色}} {\footnotesize 首席研究员兼联合创始人。推动面向化学和药物发现的量子算法研究。专注于近期量子实用性以及将量子计算与生命科学应用相连接。}
\vspace{2pt}\hrule\vspace{2pt}
\vspace{4pt}
\noindent{\bfseries Stefan Knecht}\quad{\color{graytext}\footnotesize Head of Quantum, R\&D Integration and Execution}
\par{\footnotesize\color{graytext}PhD in Theoretical Chemistry / Habilitation (Venia Legendi) in Theoretical Chemistry - Relativistic Quantum Chemistry and DMRG (2009) | Heinrich Heine University Düsseldorf, ETH Zürich \textbar\  h-index: 25-30 (in quantum chemistry; note: some aggregate }
\par{\footnotesize\color{mainred}\textbf{经历}}
\begin{itemize}[leftmargin=1.2em,itemsep=0pt,topsep=0pt]
\item {\footnotesize ETH Zürich: Postdoctoral Researcher / Habilitation, Laboratory of Physical Chemistry (Markus Reiher group) (2010-2020)}
\item {\footnotesize University of Strasbourg: Postdoctoral Researcher (2009-2010)}
\item {\footnotesize University of Southern Denmark, Odense: Postdoctoral Researcher (2008-2009)}
\item {\footnotesize Johannes Gutenberg University Mainz: Privatdozent (2020-present)}
\end{itemize}
\par{\footnotesize\color{mainred}\textbf{研究方向}}
\begin{itemize}[leftmargin=1.2em,itemsep=0pt,topsep=0pt]
\item {\footnotesize 用于量子化学的密度矩阵重整化群（DMRG）}
\item {\footnotesize 相对论量子化学与重元素}
\item {\footnotesize 变分量子本征求解器（VQE）和轨道优化方法}
\item {\footnotesize 多组态电子结构理论}
\item {\footnotesize 用于药物发现的量子嵌入方案}
\end{itemize}
\par{\footnotesize\color{mainred}\textbf{代表论文}}
\begin{itemize}[leftmargin=1.2em,itemsep=0pt,topsep=0pt]
\item {\footnotesize New Approaches for ab initio Calculations of Molecules with Strong Electron Correlation (Chimia, 2016)}
\item {\footnotesize Communication: Four-component density matrix renormalization group (J. Chem. Phys., 2014)}
\end{itemize}
\par{\footnotesize\color{mainred}\textbf{成就}}
\begin{itemize}[leftmargin=1.2em,itemsep=0pt,topsep=0pt]
\item {\footnotesize 从ETH Zürich获得理论化学Venia legendi（特许任教资格）}
\item {\footnotesize 首创了四分量相对论DMRG方法}
\item {\footnotesize 共同开发了QCMaquis DMRG软件}
\end{itemize}
\par{\footnotesize\color{mainred}\textbf{角色}} {\footnotesize 量子化学团队负责人/生命科学中量子化学问题高效量子算法开发部门主管。领导药物发现中量子化学应用的算法开发。}
\vspace{2pt}\hrule\vspace{2pt}
\vspace{4pt}
\noindent{\bfseries Daniel Cavalcanti}\quad{\color{graytext}\footnotesize Lead Quantum Information Researcher}
\par{\footnotesize\color{graytext}PhD in Theoretical Physics - Quantum Information and Quantum Correlations (2008) | ICFO - Institute of Photonic Sciences, Barcelona, Centre for Quantum Technologies, Singapore \textbar\  h-index: 35-40}
\par{\footnotesize\color{mainred}\textbf{经历}}
\begin{itemize}[leftmargin=1.2em,itemsep=0pt,topsep=0pt]
\item {\footnotesize ICFO - Institute of Photonic Sciences: Ramón y Cajal Fellow / Senior Researcher (2013-2021)}
\item {\footnotesize Centre for Quantum Technologies, Singapore: Postdoctoral Researcher / Independent Researcher (2010-2013)}
\item {\footnotesize ICFO - Institute of Photonic Sciences: Postdoctoral Researcher (2009)}
\end{itemize}
\par{\footnotesize\color{mainred}\textbf{研究方向}}
\begin{itemize}[leftmargin=1.2em,itemsep=0pt,topsep=0pt]
\item {\footnotesize 贝尔非定域性与网络中的量子关联}
\item {\footnotesize 量子导引与纠缠认证}
\item {\footnotesize 量子隐形传态与测量不相容性}
\item {\footnotesize 量子信息的半定规划}
\item {\footnotesize 用于量子态估计的张量网络方法}
\end{itemize}
\par{\footnotesize\color{mainred}\textbf{代表论文}}
\begin{itemize}[leftmargin=1.2em,itemsep=0pt,topsep=0pt]
\item {\footnotesize All Entangled States can Demonstrate Nonclassical Teleportation (Physical Review Letters, 2017)}
\item {\footnotesize Detection of entanglement in asymmetric quantum networks and multipartite quantum steering (Nature Communications, 2015)}
\end{itemize}
\par{\footnotesize\color{mainred}\textbf{成就}}
\begin{itemize}[leftmargin=1.2em,itemsep=0pt,topsep=0pt]
\item {\footnotesize 合著教科书《Semidefinite Programming in Quantum Information Science》（IOP, 2023）}
\item {\footnotesize Ramón y Cajal 著名研究基金（西班牙）}
\item {\footnotesize 108篇出版物，约6000次引用}
\end{itemize}
\par{\footnotesize\color{mainred}\textbf{角色}} {\footnotesize 首席量子信息研究员。运用量子信息理论开发药物发现算法。还设计了公司的视觉形象和品牌。}
\vspace{2pt}\hrule\vspace{2pt}
\vspace{4pt}
\noindent{\bfseries Fabijan Pavosevic}\quad{\color{graytext}\footnotesize Lead Quantum Chemist}
\par{\footnotesize\color{graytext}PhD in Theoretical Chemistry - Theoretical and Quantum Chemistry (2017) | Virginia Tech, Yale University \textbar\  h-index: 16 \textbar\  \href{https://scholar.google.com/citations?user=q2lveLEAAAAJ}{Scholar}}
\par{\footnotesize\color{mainred}\textbf{经历}}
\begin{itemize}[leftmargin=1.2em,itemsep=0pt,topsep=0pt]
\item {\footnotesize Flatiron Institute: Flatiron Research Fellow, Center for Computational Quantum Physics (CCQ) (2021-2023)}
\item {\footnotesize Yale University: Postdoctoral Associate (Sharon Hammes-Schiffer group) (2017-2021)}
\item {\footnotesize Virginia Tech: PhD Researcher (Edward Valeev group) (2012-2017)}
\item {\footnotesize University of Zagreb: BS/MS Student, Faculty of Chemical Engineering and Technology (2003-2009)}
\end{itemize}
\par{\footnotesize\color{mainred}\textbf{研究方向}}
\begin{itemize}[leftmargin=1.2em,itemsep=0pt,topsep=0pt]
\item {\footnotesize 核-电子轨道（NEO）方法}
\item {\footnotesize 量子计算的耦合簇理论}
\item {\footnotesize 幺正耦合簇（UCC）与ADAPT-VQE}
\item {\footnotesize 极化激元化学与腔修饰反应}
\item {\footnotesize 活性空间选择（AEGISS工作流）}
\end{itemize}
\par{\footnotesize\color{mainred}\textbf{代表论文}}
\begin{itemize}[leftmargin=1.2em,itemsep=0pt,topsep=0pt]
\item {\footnotesize Software for the frontiers of quantum chemistry: Q-Chem 5 (2021, 841 citations)}
\item {\footnotesize Multicomponent quantum chemistry via the NEO method (Chemical Reviews, 2020, 126 citations)}
\end{itemize}
\par{\footnotesize\color{mainred}\textbf{成就}}
\begin{itemize}[leftmargin=1.2em,itemsep=0pt,topsep=0pt]
\item {\footnotesize 计算量子物理学中心Flatiron研究所研究员}
\item {\footnotesize Q-Chem 5软件的关键贡献者（841次引用）}
\item {\footnotesize 开发了用于量子计算机上的量子化学的AEGISS活性空间选择工作流}
\end{itemize}
\par{\footnotesize\color{mainred}\textbf{角色}} {\footnotesize 首席/主力量子化学家。为近期量子设备开发量子化学算法，重点研究用于药物发现和材料科学的分子模拟。}
\vspace{2pt}\hrule\vspace{2pt}
\vspace{4pt}
\noindent{\bfseries Martina Stella}\quad{\color{graytext}\footnotesize Senior Researcher, Head of Partnerships}
\par{\footnotesize\color{graytext}PhD in Theoretical and Computational Chemistry - Theoretical and Computational Chemistry (2015) | University of Bristol, University of Cambridge \textbar\  h-index: 10 (total), 9 (recent, last 5 years) \textbar\  \href{Not publicly indexed; profile exists per AD Scientific Index}{Scholar}}
\par{\footnotesize\color{mainred}\textbf{经历}}
\begin{itemize}[leftmargin=1.2em,itemsep=0pt,topsep=0pt]
\item {\footnotesize Imperial College London: Researcher (2018–2022)}
\item {\footnotesize King's College London: Postdoctoral Researcher (2015–2018)}
\item {\footnotesize ICTP (Abdus Salam International Centre for Theoretical Physics): Boltzmann Senior Fellow / Long-term Visiting Scientist (2022–present)}
\end{itemize}
\par{\footnotesize\color{mainred}\textbf{研究方向}}
\begin{itemize}[leftmargin=1.2em,itemsep=0pt,topsep=0pt]
\item {\footnotesize 用于化学和生命科学的量子计算}
\item {\footnotesize 肿瘤光动力疗法（PDT）}
\item {\footnotesize 激发态的密度泛函理论（DFT）方法}
\item {\footnotesize BigDFT代码开发}
\item {\footnotesize 活性空间选择（AEGISS方法）}
\end{itemize}
\par{\footnotesize\color{mainred}\textbf{代表论文}}
\begin{itemize}[leftmargin=1.2em,itemsep=0pt,topsep=0pt]
\item {\footnotesize AEGISS — Atomic orbital and Entropy-based Guided Inference for Space Selection (arXiv, 2025)}
\item {\footnotesize Scalable quantum circuit generation using Majorana Propagation (ADAPT-VMPE)}
\end{itemize}
\par{\footnotesize\color{mainred}\textbf{成就}}
\begin{itemize}[leftmargin=1.2em,itemsep=0pt,topsep=0pt]
\item {\footnotesize 双学术-产业职位：在Algorithmiq全职工作的同时保持ICTP的从属关系}
\item {\footnotesize 在ICTP获得Boltzmann Fellowship独立学术职位}
\item {\footnotesize BigDFT代码的共同开发者}
\end{itemize}
\par{\footnotesize\color{mainred}\textbf{角色}} {\footnotesize 高级研究员，合作伙伴关系负责人——识别量子计算机可以模拟的生物系统，为生命科学开发NISQ软件}
\vspace{2pt}\hrule\vspace{2pt}
\vspace{4pt}
\noindent{\bfseries John Goold}\quad{\color{graytext}\footnotesize Lead Algorithmiq Ireland}
\par{\footnotesize\color{graytext}PhD in Physics (2010) - Physics — Quantum Thermodynamics (2010) | University College Cork \textbar\  h-index: 37 \textbar\  \href{https://scholar.google.com/citations?user=uPtQT_sAAAAJ}{Scholar}}
\par{\footnotesize\color{mainred}\textbf{经历}}
\begin{itemize}[leftmargin=1.2em,itemsep=0pt,topsep=0pt]
\item {\footnotesize Trinity College Dublin: Full Professor of Physics (2024–present)}
\item {\footnotesize Trinity College Dublin: Associate Professor (2019–2024)}
\item {\footnotesize Trinity College Dublin: Assistant Professor (2017–2019)}
\item {\footnotesize ICTP (Abdus Salam International Centre for Theoretical Physics): Research Scientist (2013–2017)}
\end{itemize}
\par{\footnotesize\color{mainred}\textbf{研究方向}}
\begin{itemize}[leftmargin=1.2em,itemsep=0pt,topsep=0pt]
\item {\footnotesize 量子热力学}
\item {\footnotesize 量子信息理论}
\item {\footnotesize 多体物理}
\item {\footnotesize 量子电池与充电}
\item {\footnotesize 非平衡量子引擎}
\end{itemize}
\par{\footnotesize\color{mainred}\textbf{代表论文}}
\begin{itemize}[leftmargin=1.2em,itemsep=0pt,topsep=0pt]
\item {\footnotesize The role of quantum information in thermodynamics—a topical review (2016) — 1,048 citations}
\item {\footnotesize Experimental reconstruction of work distribution and study of fluctuation relations in a closed quantum system (2014) — 432 citations}
\end{itemize}
\par{\footnotesize\color{mainred}\textbf{成就}}
\begin{itemize}[leftmargin=1.2em,itemsep=0pt,topsep=0pt]
\item {\footnotesize 与合作伙伴IBM、微软、Algorithmiq、Horizon Quantum、Moody's Analytics共同创立三一量子联盟（TQA）}
\item {\footnotesize 在都柏林圣三一学院创办爱尔兰首个量子科学与技术硕士项目}
\item {\footnotesize 当选都柏林圣三一学院院士（2022）}
\end{itemize}
\par{\footnotesize\color{mainred}\textbf{角色}} {\footnotesize Algorithmiq爱尔兰负责人——连接Algorithmiq与都柏林圣三一学院生态系统，合作使用IBM硬件进行错误缓解实验}
\vspace{2pt}\hrule\vspace{2pt}
\vspace{4pt}
\noindent{\bfseries Caterina Foti}\quad{\color{graytext}\footnotesize Senior Researcher and Outreach Coordinator}
\par{\footnotesize\color{graytext}PhD in Physics - Quantum Physics / Foundations of Quantum Mechanics (2019) | University of Florence (PhD, 2019), Aalto University (Postdoctoral Researcher) \textbar\  h-index: \textasciitilde{}8-10 (based on Nature Communications paper with \textasciitilde{} \textbar\  \href{https://www.linkedin.com/in/caterina-foti/}{LinkedIn} \textbar\  \href{https://orcid.org/0000-0002-3977-4295}{Scholar}}
\par{\footnotesize\color{mainred}\textbf{经历}}
\begin{itemize}[leftmargin=1.2em,itemsep=0pt,topsep=0pt]
\item {\footnotesize Aalto University: Postdoctoral Researcher, Department of Applied Physics (2020-present)}
\item {\footnotesize Aalto University: Intern/Visiting Researcher (during PhD) (2017)}
\item {\footnotesize Algorithmiq Ltd: Senior Researcher and Outreach Coordinator (2020-present)}
\item {\footnotesize University of Pisa: Research Collaborator ()}
\end{itemize}
\par{\footnotesize\color{mainred}\textbf{研究方向}}
\begin{itemize}[leftmargin=1.2em,itemsep=0pt,topsep=0pt]
\item {\footnotesize 量子物理教育与外展}
\item {\footnotesize 面向科学素养的量子游戏}
\item {\footnotesize 量子物理学基础（Page and Wootters mechanism）}
\item {\footnotesize 从量子纠缠中产生时间的机制}
\item {\footnotesize 文化-科学故事叙述}
\end{itemize}
\par{\footnotesize\color{mainred}\textbf{代表论文}}
\begin{itemize}[leftmargin=1.2em,itemsep=0pt,topsep=0pt]
\item {\footnotesize Time and classical equations of motion from quantum entanglement via the Page and Wootters mechanism with generalized coherent states (Nature Communic}
\item {\footnotesize Games for Quantum Physics Education (Journal of Physics: Conference Series, 2024)}
\end{itemize}
\par{\footnotesize\color{mainred}\textbf{成就}}
\begin{itemize}[leftmargin=1.2em,itemsep=0pt,topsep=0pt]
\item {\footnotesize 基于游戏和互动工具的在线量子素养平台QPlayLearn的共同创建者}
\item {\footnotesize 量子丛林互动展览（在比萨Palazzo Blu展出）的共同创建者}
\item {\footnotesize 为公众参与组织了量子游戏创作营}
\end{itemize}
\par{\footnotesize\color{mainred}\textbf{角色}} {\footnotesize 高级研究员与外展协调员 -- 负责量子外展、QPlayLearn开发、品牌推广和基础量子研究}
\vspace{2pt}\hrule\vspace{2pt}
\vspace{4pt}
\noindent{\bfseries Roberto Di Remigio Eikås}\quad{\color{graytext}\footnotesize Head of Product Engineering}
\par{\footnotesize\color{graytext}PhD - Theoretical and Computational Chemistry (Quantum Molecular Sciences) (2017) | UiT The Arctic University of Norway (Tromso), University of Pisa \textbar\  h-index: 12-13 (24+ publications, \textasciitilde{}2000+ total citations) \textbar\  \href{https://www.linkedin.com/in/robertodr/}{LinkedIn} \textbar\  \href{https://scholar.google.com.tr/citations?user=G015fgsAAAAJ}{Scholar}}
\par{\footnotesize\color{mainred}\textbf{经历}}
\begin{itemize}[leftmargin=1.2em,itemsep=0pt,topsep=0pt]
\item {\footnotesize UiT The Arctic University of Norway / Hylleraas Centre for Quantum Molecular Sciences: Postdoctoral Fellow (2017-2020)}
\item {\footnotesize Virginia Tech: Visiting Researcher (Prior to 2020)}
\item {\footnotesize University of Pisa - Molecular Modelling Lab: Researcher / Contributor (Prior to 2022)}
\item {\footnotesize Algorithmiq Ltd: Lead Software Engineer, R\&D (2022-2024)}
\end{itemize}
\par{\footnotesize\color{mainred}\textbf{研究方向}}
\begin{itemize}[leftmargin=1.2em,itemsep=0pt,topsep=0pt]
\item {\footnotesize 量子化学软件开发（PCMSolver、DIRAC、PSI4、MRChem、VAMPyR）}
\item {\footnotesize 连续溶剂化模型（极化连续介质模型）}
\item {\footnotesize 相对论量子化学（狄拉克-库仑-布莱特哈密顿量）}
\item {\footnotesize 耦合簇理论与图解蒙特卡罗方法}
\item {\footnotesize 用于电子结构的多小波方法}
\end{itemize}
\par{\footnotesize\color{mainred}\textbf{代表论文}}
\begin{itemize}[leftmargin=1.2em,itemsep=0pt,topsep=0pt]
\item {\footnotesize Full Breit Hamiltonian in the Multiwavelets Framework - J. Chem. Theory Comput. (2024)}
\item {\footnotesize VAMPyR - A high-level Python library for mathematical operations in a multiwavelet representation - J. Chem. Phys. (2024)}
\end{itemize}
\par{\footnotesize\color{mainred}\textbf{成就}}
\begin{itemize}[leftmargin=1.2em,itemsep=0pt,topsep=0pt]
\item {\footnotesize PCMSolver（溶剂化建模开源库）的主要开发者}
\item {\footnotesize 多个旗舰量子化学程序（PSI4、DIRAC、Dalton Project）的重要贡献者}
\item {\footnotesize 科学计算领域广泛使用的《CMake Cookbook》合著者}
\end{itemize}
\par{\footnotesize\color{mainred}\textbf{角色}} {\footnotesize 产品工程负责人（曾任首席软件工程师，研发岗）——连接量子化学研究与产品开发，主导Algorithmiq量子计算平台的软件工程}
\vspace{2pt}\hrule\vspace{2pt}
\vspace{4pt}
\noindent{\bfseries Kirsten Nehr}\quad{\color{graytext}\footnotesize Chief Operating Officer}
\par{\footnotesize\color{graytext}Master's degree - Biological Sciences (Not publicly specified) | University of Oxford \textbar\  h-index: N/A \textbar\  \href{https://www.linkedin.com/in/kirstennehr/}{LinkedIn} \textbar\  \href{N/A}{Scholar}}
\par{\footnotesize\color{mainred}\textbf{经历}}
\begin{itemize}[leftmargin=1.2em,itemsep=0pt,topsep=0pt]
\item {\footnotesize Healthcare and Technology sectors: Commercial roles (marketing, branding, life sciences partnerships, business strategy, operations) (10+ years)}
\end{itemize}
\par{\footnotesize\color{mainred}\textbf{成就}}
\begin{itemize}[leftmargin=1.2em,itemsep=0pt,topsep=0pt]
\item {\footnotesize 在AI for Good（国际电信联盟）平台发表演讲}
\item {\footnotesize 在Algorithmiq打造定义行业类别的量子计算业务}
\end{itemize}
\par{\footnotesize\color{mainred}\textbf{角色}} {\footnotesize 首席运营官（COO）/ 幕僚长——在专注于医疗/生命科学的量子计算与人工智能公司中，负责运营、业务战略、市场营销、品牌建设及生命科学合作伙伴关系}
\vspace{2pt}\hrule\vspace{2pt}
\vspace{4pt}
\noindent{\bfseries Ben Bakker}\quad{\color{graytext}\footnotesize Chief Financial Officer}
\par{\footnotesize\color{graytext}Not publicly available - Not publicly available (likely Finance / Business) (Not publicly available) \textbar\  h-index: N/A \textbar\  \href{https://www.linkedin.com/in/benjaminbakkernz/}{LinkedIn} \textbar\  \href{N/A}{Scholar}}
\par{\footnotesize\color{mainred}\textbf{成就}}
\begin{itemize}[leftmargin=1.2em,itemsep=0pt,topsep=0pt]
\item {\footnotesize 量子计算初创公司Algorithmiq的首席财务官}
\end{itemize}
\par{\footnotesize\color{mainred}\textbf{角色}} {\footnotesize 首席财务官 - 负责Algorithmiq的财务战略、融资及财务运营}
\vspace{2pt}\hrule\vspace{2pt}
\vspace{4pt}
\noindent{\bfseries Irene Gentili}\quad{\color{graytext}\footnotesize Chief of Staff}
\par{\footnotesize\color{graytext}Master's degree - Business / Management (Not publicly specified) | Bocconi University (Milan) - Bachelor's, ESADE Business School (Barcelona) - Master's \textbar\  h-index: N/A \textbar\  \href{https://www.linkedin.com/in/irenegentili/}{LinkedIn} \textbar\  \href{N/A}{Scholar}}
\par{\footnotesize\color{mainred}\textbf{经历}}
\begin{itemize}[leftmargin=1.2em,itemsep=0pt,topsep=0pt]
\item {\footnotesize Redstone: VC Investor (\textasciitilde{}2.5 years)}
\item {\footnotesize Speedinvest: Associate, Marketplaces \& Consumer team ()}
\item {\footnotesize Samaipata: VC Investor ()}
\item {\footnotesize Consulting (Milan): Consultant (Early career)}
\end{itemize}
\par{\footnotesize\color{mainred}\textbf{成就}}
\begin{itemize}[leftmargin=1.2em,itemsep=0pt,topsep=0pt]
\item {\footnotesize 登上《Woman Taught Tech》播客，探讨风险投资领域的多样性和对女性创始人的资金支持}
\item {\footnotesize 从风险投资领域转型至深科技量子初创企业的运营领导层}
\end{itemize}
\par{\footnotesize\color{mainred}\textbf{角色}} {\footnotesize 幕僚长 - 负责构建运营架构、流程和日常运营，以支持公司规模化发展阶段。常驻米兰（Algorithmiq全球总部）。}
\vspace{2pt}\hrule\vspace{2pt}
\vspace{4pt}
\noindent{\bfseries Minna Gunes}\quad{\color{graytext}\footnotesize Head of People and Strategic Funding}
\par{\footnotesize\color{graytext}PhD in Natural Sciences - Natural Sciences (2004) | University of Jyvaskyla - PhD (2004) \textbar\  h-index: N/A (career focused on academic coordination rathe \textbar\  \href{https://www.linkedin.com/in/minna-gunes/}{LinkedIn} \textbar\  \href{N/A (limited publication record)}{Scholar}}
\par{\footnotesize\color{mainred}\textbf{经历}}
\begin{itemize}[leftmargin=1.2em,itemsep=0pt,topsep=0pt]
\item {\footnotesize Aalto University: Academic Coordinator for InstituteQ, Quantum Technology Finland, OtaNano, Centre for Quantum Engineering (\textasciitilde{}2015-2024)}
\item {\footnotesize University of Helsinki: Researcher/Coordinator ()}
\end{itemize}
\par{\footnotesize\color{mainred}\textbf{研究方向}}
\begin{itemize}[leftmargin=1.2em,itemsep=0pt,topsep=0pt]
\item {\footnotesize 量子生态系统协调与基础设施管理}
\end{itemize}
\par{\footnotesize\color{mainred}\textbf{代表论文}}
\begin{itemize}[leftmargin=1.2em,itemsep=0pt,topsep=0pt]
\item {\footnotesize InstituteQ - The Finnish National Quantum Institute - Arkhimedes (2022), co-authored with J.P. Pekola, H.S. Majumdar, P. Muratore Ginanneschi}
\end{itemize}
\par{\footnotesize\color{mainred}\textbf{成就}}
\begin{itemize}[leftmargin=1.2em,itemsep=0pt,topsep=0pt]
\item {\footnotesize 芬兰国家量子生态系统计划（InstituteQ、Quantum Technology Finland、OtaNano）的关键协调人}
\item {\footnotesize 多语能力：芬兰语（母语），英语（流利），瑞典语（一般），德语/西班牙语（基础），土耳其语（良好口语，书面语一般）}
\item {\footnotesize 代表Algorithmiq参与北欧-波罗的海量子生态系统报告（Nordic Innovation）}
\end{itemize}
\par{\footnotesize\color{mainred}\textbf{角色}} {\footnotesize 人才与战略资助负责人（在Algorithmiq团队页面上也同时列为战略联盟负责人兼人才合伙人）——负责人才运营、战略伙伴关系和资助协调}
\vspace{2pt}\hrule\vspace{2pt}
\vspace{4pt}
\noindent{\bfseries Emilia Silius}\quad{\color{graytext}\footnotesize Business Operations Coordinator}
\par{\footnotesize\color{graytext}Master's degree - Marketing / Business (2024) | Aalto University School of Business - Master's (2024), Aalto University School of Business - Bachelor's in Marketing (2021) \textbar\  h-index: N/A \textbar\  \href{https://www.linkedin.com/in/emilia-silius/}{LinkedIn} \textbar\  \href{N/A}{Scholar}}
\par{\footnotesize\color{mainred}\textbf{经历}}
\begin{itemize}[leftmargin=1.2em,itemsep=0pt,topsep=0pt]
\item {\footnotesize Aalto University: Student (Bachelor's and Master's theses) (2021-2024)}
\end{itemize}
\par{\footnotesize\color{mainred}\textbf{研究方向}}
\begin{itemize}[leftmargin=1.2em,itemsep=0pt,topsep=0pt]
\item {\footnotesize Consumer behavior and retail therapy (Bachelor's thesis)}
\item {\footnotesize Slow fashion and sustainable consumption (Master's thesis)}
\end{itemize}
\par{\footnotesize\color{mainred}\textbf{代表论文}}
\begin{itemize}[leftmargin=1.2em,itemsep=0pt,topsep=0pt]
\item {\footnotesize Bachelor's thesis: 'Can Money Buy Happiness?' - A literature review on retail therapy (2021)}
\item {\footnotesize Master's thesis: 'Fashion is temporary, but style endures: Exploring the slow fashion phenomenon and efforts of sustainable clothing brands toward mor}
\end{itemize}
\par{\footnotesize\color{mainred}\textbf{成就}}
\begin{itemize}[leftmargin=1.2em,itemsep=0pt,topsep=0pt]
\item {\footnotesize Completed both Bachelor's and Master's degrees at Aalto University School of Business}
\item {\footnotesize Master's thesis involved qualitative research with Nordic sustainable fashion brands}
\end{itemize}
\par{\footnotesize\color{mainred}\textbf{角色}} {\footnotesize Business Operations Coordinator - responsible for business operations support at the Helsinki office}
\vspace{2pt}\hrule\vspace{2pt}
\subsection{研发团队（19人）}
\vspace{4pt}
\noindent{\bfseries Tommaso Tufarelli}\quad{\color{graytext}\footnotesize Senior Researcher and R\&D Project Manager}
\par{\footnotesize\color{graytext}PhD (Theoretical Quantum Physics) - Quantum Optics and Open Quantum Systems () | University of Nottingham, Imperial College London \textbar\  h-index: Not available — likely moderate based on publicati \textbar\  \href{Not found in public search}{Scholar}}
\par{\footnotesize\color{mainred}\textbf{经历}}
\begin{itemize}[leftmargin=1.2em,itemsep=0pt,topsep=0pt]
\item {\footnotesize University of Nottingham: Assistant Professor, School of Mathematical Sciences (\textasciitilde{}2021–present)}
\item {\footnotesize Imperial College London: Researcher, QOLS Blackett Laboratory (–)}
\item {\footnotesize University College London: Researcher, Department of Physics \& Astronomy (–)}
\item {\footnotesize Universitat Ulm: Researcher, Institut fur Theoretische Physik (–)}
\end{itemize}
\par{\footnotesize\color{mainred}\textbf{研究方向}}
\begin{itemize}[leftmargin=1.2em,itemsep=0pt,topsep=0pt]
\item {\footnotesize 量子光学与开放量子系统}
\item {\footnotesize 非马尔可夫量子动力学}
\item {\footnotesize 量子计量学与量子费希尔信息}
\item {\footnotesize 量子关联与量子失谐}
\item {\footnotesize 腔量子电动力学与电路量子电动力学}
\end{itemize}
\par{\footnotesize\color{mainred}\textbf{代表论文}}
\begin{itemize}[leftmargin=1.2em,itemsep=0pt,topsep=0pt]
\item {\footnotesize General Expressions for the Quantum Fisher Information Matrix with Applications to Discrete Quantum Imaging (PRX Quantum, 2021)}
\item {\footnotesize Coherently opening a high-Q cavity (Physical Review Letters, 2013)}
\end{itemize}
\par{\footnotesize\color{mainred}\textbf{成就}}
\begin{itemize}[leftmargin=1.2em,itemsep=0pt,topsep=0pt]
\item {\footnotesize 因量子关联与不确定性方面的工作登上Physical Review Letters（2013年）封面}
\item {\footnotesize 量子计量学和费希尔信息理论专家}
\end{itemize}
\par{\footnotesize\color{mainred}\textbf{角色}} {\footnotesize 高级研究员兼研发项目经理（用户所述职务；公开网络无直接确认其在Algorithmiq任职的信息——公开资料显示其为诺丁汉大学助理教授）}
\vspace{2pt}\hrule\vspace{2pt}
\vspace{4pt}
\noindent{\bfseries Sergei Filippov}\quad{\color{graytext}\footnotesize Senior Researcher}
\par{\footnotesize\color{graytext}PhD in Theoretical Physics (2012) - Theoretical Physics — Quantum Information Theory (2012) | Moscow Institute of Physics and Technology (MIPT) \textbar\  h-index: 15–20 (estimated; \textasciitilde{}70 papers, \textasciitilde{}1,000 citations on  \textbar\  \href{https://scholar.google.com/citations?user=AimFBr8AAAAJ&hl=en}{Scholar}}
\par{\footnotesize\color{mainred}\textbf{经历}}
\begin{itemize}[leftmargin=1.2em,itemsep=0pt,topsep=0pt]
\item {\footnotesize Moscow Institute of Physics and Technology (MIPT): Associate Professor (–)}
\item {\footnotesize Steklov Mathematical Institute, Russian Academy of Sciences: Senior Researcher (2018–)}
\end{itemize}
\par{\footnotesize\color{mainred}\textbf{研究方向}}
\begin{itemize}[leftmargin=1.2em,itemsep=0pt,topsep=0pt]
\item {\footnotesize Open quantum systems and non-Markovian dynamics}
\item {\footnotesize Quantum error mitigation via tensor networks}
\item {\footnotesize Quantum measurements and tomography}
\item {\footnotesize Quantum entanglement and quantum channels}
\item {\footnotesize Quantum communication (coherent information superadditivity)}
\end{itemize}
\par{\footnotesize\color{mainred}\textbf{代表论文}}
\begin{itemize}[leftmargin=1.2em,itemsep=0pt,topsep=0pt]
\item {\footnotesize Scalable tensor-network error mitigation for near-term quantum computing (arXiv:2307.11740) — foundational TEM method}
\item {\footnotesize Machine learning non-Markovian quantum dynamics (Physical Review Letters, 2020) — 79 citations}
\end{itemize}
\par{\footnotesize\color{mainred}\textbf{成就}}
\begin{itemize}[leftmargin=1.2em,itemsep=0pt,topsep=0pt]
\item {\footnotesize Core inventor of Algorithmiq's Tensor-network Error Mitigation (TEM) method, commercially launched on IBM Qiskit Functions Catalog (Sept 2024)}
\item {\footnotesize Co-author with Alexander S. Holevo on quantum Gaussian maximizers}
\item {\footnotesize Born June 19, 1987}
\end{itemize}
\par{\footnotesize\color{mainred}\textbf{角色}} {\footnotesize Senior Researcher — leads tensor network noise characterization and error mitigation research}
\vspace{2pt}\hrule\vspace{2pt}
\vspace{4pt}
\noindent{\bfseries Adam Glos}\quad{\color{graytext}\footnotesize Computer Science Researcher}
\par{\footnotesize\color{graytext}PhD in Computer Science - Computer Science — Quantum Walks and Quantum Algorithms () | Institute of Theoretical and Applied Informatics, Polish Academy of Sciences, Silesian University of Technology \textbar\  h-index: 12–15 (estimated from Google Scholar citation coun \textbar\  \href{https://scholar.google.com/citations?user=3f3cR0AAAAAJ}{Scholar}}
\par{\footnotesize\color{mainred}\textbf{经历}}
\begin{itemize}[leftmargin=1.2em,itemsep=0pt,topsep=0pt]
\item {\footnotesize Institute of Theoretical and Applied Informatics, Polish Academy of Sciences: PhD Researcher (–)}
\item {\footnotesize QWorld Association: QResearch Coordinator (–)}
\end{itemize}
\par{\footnotesize\color{mainred}\textbf{研究方向}}
\begin{itemize}[leftmargin=1.2em,itemsep=0pt,topsep=0pt]
\item {\footnotesize 图上的量子行走与空间搜索}
\item {\footnotesize 量子优化（QAOA，量子退火）}
\item {\footnotesize 变分量子算法}
\item {\footnotesize 量子电路编译与态制备}
\item {\footnotesize 量子网络医学}
\end{itemize}
\par{\footnotesize\color{mainred}\textbf{代表论文}}
\begin{itemize}[leftmargin=1.2em,itemsep=0pt,topsep=0pt]
\item {\footnotesize Prog-QAOA: Framework for resource-efficient quantum optimization through classical programs (Quantum 9:1663, 2025)}
\item {\footnotesize Treespilation: Architecture- and State-Optimised Fermion-to-Qubit Mappings}
\end{itemize}
\par{\footnotesize\color{mainred}\textbf{成就}}
\begin{itemize}[leftmargin=1.2em,itemsep=0pt,topsep=0pt]
\item {\footnotesize 与Ozlem Salehi Koken在Classiq编码挑战赛（2022）中获得金牌}
\item {\footnotesize Algorithmiq公司态制备团队负责人}
\item {\footnotesize 在arXiv上发表了50多篇论文}
\end{itemize}
\par{\footnotesize\color{mainred}\textbf{角色}} {\footnotesize 计算机科学研究员 — 态制备团队负责人，专注于NISQ设备的电路编译与优化}
\vspace{2pt}\hrule\vspace{2pt}
\vspace{4pt}
\noindent{\bfseries Walter Talarico}\quad{\color{graytext}\footnotesize Senior Researcher}
\par{\footnotesize\color{graytext}PhD in Theoretical Physics (2020) - Theoretical Physics — Condensed Matter and Quantum Many-Body Physics (2020) | University of Turku, Universita degli Studi della Calabria \textbar\  h-index: 3–6 (estimated from partial citation data) \textbar\  \href{https://scholar.google.co.th/citations?user=I9J5N84AAAAJ}{Scholar}}
\par{\footnotesize\color{mainred}\textbf{经历}}
\begin{itemize}[leftmargin=1.2em,itemsep=0pt,topsep=0pt]
\item {\footnotesize Aalto University: Postdoctoral Researcher, Department of Applied Physics (2021–2023)}
\item {\footnotesize University of Turku: PhD Student, Department of Physics and Astronomy (2017–2020)}
\end{itemize}
\par{\footnotesize\color{mainred}\textbf{研究方向}}
\begin{itemize}[leftmargin=1.2em,itemsep=0pt,topsep=0pt]
\item {\footnotesize 化学领域的量子算法（VQE、ADAPT-VQE、qEOM）}
\item {\footnotesize 量子多体物理与开放量子系统}
\item {\footnotesize 非平衡格林函数}
\item {\footnotesize 纳米尺度系统中的量子输运}
\item {\footnotesize 变分量子本征求解器方法}
\end{itemize}
\par{\footnotesize\color{mainred}\textbf{代表论文}}
\begin{itemize}[leftmargin=1.2em,itemsep=0pt,topsep=0pt]
\item {\footnotesize Toward Accurate Post-Born–Oppenheimer Molecular Simulations on Quantum Computers (J. Chem. Theory Comput., 2023)}
\item {\footnotesize Self-Consistent Field Approach for the Variational Quantum Eigensolver: Orbital Optimization Goes Adaptive (J. Phys. Chem. A, 2024)}
\end{itemize}
\par{\footnotesize\color{mainred}\textbf{成就}}
\begin{itemize}[leftmargin=1.2em,itemsep=0pt,topsep=0pt]
\item {\footnotesize 开发了具有可极化嵌入的VQE-SCF方法，用于真实化学体系}
\item {\footnotesize 在量子中心的强关联和动态电子关联方面做出贡献（面向近期设备的NEVPT2）}
\item {\footnotesize 开放量子系统和非平衡格林函数专家}
\end{itemize}
\par{\footnotesize\color{mainred}\textbf{角色}} {\footnotesize 高级研究员 — 为近期量子计算机设计优化软件和算法，专注于量子化学应用}
\vspace{2pt}\hrule\vspace{2pt}
\vspace{4pt}
\noindent{\bfseries Zoltán Zimborás}\quad{\color{graytext}\footnotesize Senior Researcher}
\par{\footnotesize\color{graytext}PhD - Quantum Information Science () | Eotvos University, Budapest, UPV Bilbao \textbar\  h-index: 15–20 (estimated from publication citation counts) \textbar\  \href{Not directly found via public search (likely exists but URL not indexed)}{Scholar}}
\par{\footnotesize\color{mainred}\textbf{经历}}
\begin{itemize}[leftmargin=1.2em,itemsep=0pt,topsep=0pt]
\item {\footnotesize Wigner Research Centre for Physics, Budapest: Leader, Quantum Computing and Information Research Group (–)}
\item {\footnotesize Eotvos Lorand University: Faculty of Informatics (–)}
\item {\footnotesize University of Helsinki: Researcher, Department of Physics (–)}
\item {\footnotesize UPV Bilbao: Postdoctoral Researcher (–)}
\end{itemize}
\par{\footnotesize\color{mainred}\textbf{研究方向}}
\begin{itemize}[leftmargin=1.2em,itemsep=0pt,topsep=0pt]
\item {\footnotesize 量子门分解与编译}
\item {\footnotesize 量子优势方案与计算普适性}
\item {\footnotesize 变分量子算法}
\item {\footnotesize 量子误差缓解（读出噪声）}
\item {\footnotesize 分子中的量子关联}
\end{itemize}
\par{\footnotesize\color{mainred}\textbf{代表论文}}
\begin{itemize}[leftmargin=1.2em,itemsep=0pt,topsep=0pt]
\item {\footnotesize Approaching the theoretical limit in quantum gate decomposition (Quantum 6:710, 2022)}
\item {\footnotesize Mitigation of readout noise in near-term quantum devices by classical post-processing based on detector tomography (Quantum 4:257, 2020)}
\end{itemize}
\par{\footnotesize\color{mainred}\textbf{成就}}
\begin{itemize}[leftmargin=1.2em,itemsep=0pt,topsep=0pt]
\item {\footnotesize QWorld（非营利性量子编程教育组织）董事会成员}
\item {\footnotesize NKFI资助项目“近期量子计算机”首席研究员（2020–2025，约3550万匈牙利福林）}
\item {\footnotesize 匈牙利量子信息国家实验室（QNL）成员}
\end{itemize}
\par{\footnotesize\color{mainred}\textbf{角色}} {\footnotesize 高级研究员——量子算法设计、误差缓解与门分解}
\vspace{2pt}\hrule\vspace{2pt}
\vspace{4pt}
\noindent{\bfseries Pi Haase}\quad{\color{graytext}\footnotesize Computational Chemistry Researcher}
\par{\footnotesize\color{graytext}PhD in Chemistry - Relativistic Quantum Chemistry / Electronic Structure Theory (2021) | University of Groningen (PhD, 2021), University of Copenhagen (BSc/MSc, likely) \textbar\  h-index: Not publicly available (early career, \textasciitilde{}15 publicat \textbar\  \href{https://www.linkedin.com/in/pi-a-b-haase/}{LinkedIn} \textbar\  \href{https://arxivlens.com/Author/Details/pi-a.-b.-haase-9162-5883937}{Scholar}}
\par{\footnotesize\color{mainred}\textbf{经历}}
\begin{itemize}[leftmargin=1.2em,itemsep=0pt,topsep=0pt]
\item {\footnotesize Cleveland State University: Visiting/Research Affiliate (2024-2025)}
\item {\footnotesize University of Groningen: PhD Researcher (2017-2021)}
\item {\footnotesize UiT The Arctic University of Norway: Research Collaborator (prior to 2021)}
\end{itemize}
\par{\footnotesize\color{mainred}\textbf{研究方向}}
\begin{itemize}[leftmargin=1.2em,itemsep=0pt,topsep=0pt]
\item {\footnotesize 相对论耦合簇方法}
\item {\footnotesize 超精细结构常数与不确定度估计}
\item {\footnotesize 分子宇称不守恒}
\item {\footnotesize 多参考量子化学的活性空间选择（AEGISS方法）}
\item {\footnotesize 量子计算在化学中的应用}
\end{itemize}
\par{\footnotesize\color{mainred}\textbf{代表论文}}
\begin{itemize}[leftmargin=1.2em,itemsep=0pt,topsep=0pt]
\item {\footnotesize AEGISS -- Atomic orbital and Entropy-based Guided Inference for Space Selection (arXiv, 2025)}
\item {\footnotesize Hyperfine Structure Constants on the Relativistic Coupled Cluster Level with Associated Uncertainties (J. Phys. Chem. A, 2020)}
\end{itemize}
\par{\footnotesize\color{mainred}\textbf{成就}}
\begin{itemize}[leftmargin=1.2em,itemsep=0pt,topsep=0pt]
\item {\footnotesize 开发了AEGISS方法，用于半自动活性空间选择，将轨道熵分析与原子轨道投影统一起来}
\item {\footnotesize 提出分子性质计算的第一性原理不确定度估计方案，可实现约5.5\%的保守不确定度界限}
\item {\footnotesize 架起量子化学与药物发现中量子计算应用之间的桥梁}
\end{itemize}
\par{\footnotesize\color{mainred}\textbf{角色}} {\footnotesize 计算化学研究员——开发用于分子模拟中量子计算应用的量子化学方法}
\vspace{2pt}\hrule\vspace{2pt}
\vspace{4pt}
\noindent{\bfseries Özlem Salehi}\quad{\color{graytext}\footnotesize Computer Science Researcher}
\par{\footnotesize\color{graytext}PhD in Theoretical Computer Science - Computer Engineering / Theoretical Computer Science (2019) | Boğaziçi University, Istanbul (PhD, 2019), Boğaziçi University, Istanbul (MSc, 2013) \textbar\  h-index: \textasciitilde{}5 (384 total citations per AD Scientific Index) \textbar\  \href{https://www.linkedin.com/in/ozlem-salehi/}{LinkedIn} \textbar\  \href{https://ieeexplore.ieee.org/author/37088516883}{Scholar}}
\par{\footnotesize\color{mainred}\textbf{经历}}
\begin{itemize}[leftmargin=1.2em,itemsep=0pt,topsep=0pt]
\item {\footnotesize Institute of Theoretical and Applied Informatics, Polish Academy of Sciences (IITiS PAN): Postdoctoral Researcher (2021-2023)}
\item {\footnotesize Özyeğin University: Instructor, Department of Computer Science (2019-2021)}
\item {\footnotesize Algorithmiq: Computer Science Researcher (2023-present)}
\end{itemize}
\par{\footnotesize\color{mainred}\textbf{研究方向}}
\begin{itemize}[leftmargin=1.2em,itemsep=0pt,topsep=0pt]
\item {\footnotesize 量子优化（QAOA、量子退火、VQE）}
\item {\footnotesize 用于组合优化的QUBO和HOBO公式化}
\item {\footnotesize 量子计算教育}
\item {\footnotesize 量子有限自动机与计算复杂性}
\end{itemize}
\par{\footnotesize\color{mainred}\textbf{代表论文}}
\begin{itemize}[leftmargin=1.2em,itemsep=0pt,topsep=0pt]
\item {\footnotesize Prog-QAOA: Framework for resource-efficient quantum optimization through classical programs (Quantum, Vol. 9, 2025)}
\item {\footnotesize Optimizing the Production of Test Vehicles Using Hybrid Constrained Quantum Annealing (SN Computer Science, 2023)}
\end{itemize}
\par{\footnotesize\color{mainred}\textbf{成就}}
\begin{itemize}[leftmargin=1.2em,itemsep=0pt,topsep=0pt]
\item {\footnotesize QTurkey联合创始人，该社区旨在提升土耳其的量子技术意识}
\item {\footnotesize 组织量子编程工作坊并开发开源教育材料}
\item {\footnotesize 最高引用量子优化论文被引用51次}
\end{itemize}
\par{\footnotesize\color{mainred}\textbf{角色}} {\footnotesize 计算机科学研究员——开发用于生命科学应用的量子优化算法}
\vspace{2pt}\hrule\vspace{2pt}
\vspace{4pt}
\noindent{\bfseries Alessio Fallani}\quad{\color{graytext}\footnotesize Researcher}
\par{\footnotesize\color{graytext}PhD in Physics - Quantum Physics / Machine Learning for Drug Discovery (2025) | University of Luxembourg (PhD, 2025), University of Milan, Dipartimento di Fisica Aldo Pontremoli \textbar\  h-index: \textasciitilde{}4 (early career, Nature Communications paper) \textbar\  \href{https://www.linkedin.com/in/alessiofallani/}{LinkedIn} \textbar\  \href{https://orbilu.uni.lu/profile?uid=50045086}{Scholar}}
\par{\footnotesize\color{mainred}\textbf{经历}}
\begin{itemize}[leftmargin=1.2em,itemsep=0pt,topsep=0pt]
\item {\footnotesize Algorithmiq Ltd: Researcher (2024-present)}
\item {\footnotesize University of Luxembourg: PhD Researcher (supervisor: Prof. Alexandre Tkatchenko) (2021-2025)}
\item {\footnotesize University of Milan: Research Affiliate (prior to 2021)}
\end{itemize}
\par{\footnotesize\color{mainred}\textbf{研究方向}}
\begin{itemize}[leftmargin=1.2em,itemsep=0pt,topsep=0pt]
\item {\footnotesize 用于药物发现的量子化学和机器学习}
\item {\footnotesize 用于分子设计的量子逆映射（QIM）}
\item {\footnotesize 在量子性质上预训练的图变换器用于ADMET建模}
\item {\footnotesize 用于过渡金属光敏剂发现的主动学习}
\item {\footnotesize 光子量子计算模拟非共价相互作用}
\end{itemize}
\par{\footnotesize\color{mainred}\textbf{代表论文}}
\begin{itemize}[leftmargin=1.2em,itemsep=0pt,topsep=0pt]
\item {\footnotesize Inverse Mapping of Quantum Properties to Structures for Chemical Space of Small Organic Molecules (Nature Communications, 2024)}
\item {\footnotesize Pretraining Graph Transformers with Atom-in-a-Molecule Quantum Properties for Improved ADMET Modeling (Journal of Cheminformatics, 2025)}
\end{itemize}
\par{\footnotesize\color{mainred}\textbf{成就}}
\begin{itemize}[leftmargin=1.2em,itemsep=0pt,topsep=0pt]
\item {\footnotesize 关于量子逆映射分子设计的 Nature Communications 论文}
\item {\footnotesize 由Janssen Pharma联合指导的博士研究（产业合作）}
\item {\footnotesize 开发了用于类药分子量子力学探索的Aquamarine（AQM）数据集}
\end{itemize}
\par{\footnotesize\color{mainred}\textbf{角色}} {\footnotesize 研究员——将量子化学和机器学习应用于药物发现}
\vspace{2pt}\hrule\vspace{2pt}
\vspace{4pt}
\noindent{\bfseries Fabian Langkabel}\quad{\color{graytext}\footnotesize Researcher}
\par{\footnotesize\color{graytext}Dr. rer. nat. (PhD) in Chemistry - Theoretical Chemistry / Quantum Computing (2023) | Freie Universität Berlin (PhD, 2023), Helmholtz-Zentrum Berlin (PhD research) \textbar\  h-index: \textasciitilde{}4-5 (early career) \textbar\  \href{https://www.linkedin.com/in/fabian-langkabel-728bba275/}{LinkedIn} \textbar\  \href{Not publicly found; indexed via Semantic Scholar and arXiv}{Scholar}}
\par{\footnotesize\color{mainred}\textbf{经历}}
\begin{itemize}[leftmargin=1.2em,itemsep=0pt,topsep=0pt]
\item {\footnotesize Algorithmiq Ltd: Researcher (2024-present)}
\item {\footnotesize Helmholtz-Zentrum Berlin (HZB): PhD Researcher (HEIBRiDs fellowship) (2020-2023)}
\item {\footnotesize Freie Universität Berlin: Doctoral Candidate, Institute of Chemistry and Biochemistry (2020-2023)}
\end{itemize}
\par{\footnotesize\color{mainred}\textbf{研究方向}}
\begin{itemize}[leftmargin=1.2em,itemsep=0pt,topsep=0pt]
\item {\footnotesize 分子中电子动力学的量子算法}
\item {\footnotesize 多分辨率变分量子本征值求解器（VQE）}
\item {\footnotesize 激光驱动的电子动力学模拟}
\item {\footnotesize Jellyfish模块化模拟代码开发}
\item {\footnotesize 粒子间库仑衰变（ICD）动力学}
\end{itemize}
\par{\footnotesize\color{mainred}\textbf{代表论文}}
\begin{itemize}[leftmargin=1.2em,itemsep=0pt,topsep=0pt]
\item {\footnotesize The advent of fully variational quantum eigensolvers using a hybrid multiresolution approach (arXiv, 2024)}
\item {\footnotesize Jellyfish: A modular code for wave function-based electron dynamics simulations and visualizations on traditional and quantum compute architectures (W}
\end{itemize}
\par{\footnotesize\color{mainred}\textbf{成就}}
\begin{itemize}[leftmargin=1.2em,itemsep=0pt,topsep=0pt]
\item {\footnotesize 首次实现分子中精确激光驱动电子动力学的量子计算算法}
\item {\footnotesize 开发了支持经典与量子架构的Jellyfish模拟代码}
\item {\footnotesize 入选HZB新闻稿（2022年11月）“量子算法节省电子动力学计算时间”}
\end{itemize}
\par{\footnotesize\color{mainred}\textbf{角色}} {\footnotesize 研究员 -- 将量子算法应用于计算化学，以实现药物研发中的分子模拟}
\vspace{2pt}\hrule\vspace{2pt}
\vspace{4pt}
\noindent{\bfseries Anton Nykänen}\quad{\color{graytext}\footnotesize Researcher}
\par{\footnotesize\color{graytext}MSc in Theoretical Physics - Theoretical Physics / Quantum Computing (2021) | University of Turku (MSc, 2021), University of Turku (BSc) \textbar\  h-index: \textasciitilde{}4-5 (early career, multiple JPC/JCTC publications \textbar\  \href{https://www.linkedin.com/in/anton-nykänen-6b434a203/}{LinkedIn} \textbar\  \href{https://orcid.org/0009-0003-1949-3028}{Scholar}}
\par{\footnotesize\color{mainred}\textbf{经历}}
\begin{itemize}[leftmargin=1.2em,itemsep=0pt,topsep=0pt]
\item {\footnotesize Algorithmiq Oy: Junior Researcher / Researcher (2021-present)}
\end{itemize}
\par{\footnotesize\color{mainred}\textbf{研究方向}}
\begin{itemize}[leftmargin=1.2em,itemsep=0pt,topsep=0pt]
\item {\footnotesize 变分量子本征求解器（VQE）和ADAPT-VQE方法}
\item {\footnotesize 量子计算机上分子的激发能计算}
\item {\footnotesize 后玻恩-奥本海默分子模拟}
\item {\footnotesize 费米子到量子位的映射与马约拉纳传播}
\item {\footnotesize 量子错误缓解策略}
\end{itemize}
\par{\footnotesize\color{mainred}\textbf{代表论文}}
\begin{itemize}[leftmargin=1.2em,itemsep=0pt,topsep=0pt]
\item {\footnotesize Scalable quantum circuit generation for iterative ground state approximation using Majorana Propagation (arXiv, 2026)}
\item {\footnotesize Simulation of Fermionic circuits using Majorana Propagation (arXiv, 2025)}
\end{itemize}
\par{\footnotesize\color{mainred}\textbf{成就}}
\begin{itemize}[leftmargin=1.2em,itemsep=0pt,topsep=0pt]
\item {\footnotesize Algorithmiq的ADAPT-VQE方法论的核心贡献者}
\item {\footnotesize 与AstraZeneca、IBM Research、Trinity College Dublin和Cleveland Clinic的合作}
\item {\footnotesize 硕士论文：利用近期量子计算机探索纠缠现象（University of Turku，2021年）}
\end{itemize}
\par{\footnotesize\color{mainred}\textbf{角色}} {\footnotesize Algorithmiq的研究员（初级研究员）——开发用于药物发现中分子模拟的量子算法}
\vspace{2pt}\hrule\vspace{2pt}
\vspace{4pt}
\noindent{\bfseries Keijo Korhonen}\quad{\color{graytext}\footnotesize Researcher}
\par{\footnotesize\color{graytext}MSc in Theoretical and Computational Methods - Quantum Information / Theoretical Physics (ongoing PhD (started 2023)) | University of Helsinki (PhD ongoing, 2023-present), University of Helsinki (MSc, 2022) \textbar\  h-index: \textasciitilde{}3-4 (early career, npj Quantum Information public \textbar\  \href{https://www.linkedin.com/in/keijokorhonen/}{LinkedIn} \textbar\  \href{https://orcid.org/0009-0003-2647-3105}{Scholar}}
\par{\footnotesize\color{mainred}\textbf{经历}}
\begin{itemize}[leftmargin=1.2em,itemsep=0pt,topsep=0pt]
\item {\footnotesize Algorithmiq Ltd: Junior Researcher (2022-present)}
\item {\footnotesize University of Helsinki: PhD Student, QTF Centre of Excellence (2023-present)}
\end{itemize}
\par{\footnotesize\color{mainred}\textbf{研究方向}}
\begin{itemize}[leftmargin=1.2em,itemsep=0pt,topsep=0pt]
\item {\footnotesize 近期量子硬件上的高精度测量}
\item {\footnotesize 随机测量与经典影子}
\item {\footnotesize 量子测量层析（QMT）}
\item {\footnotesize 基于局部最优对偶框架的影子估计}
\item {\footnotesize 分子能量估计（VQE）}
\end{itemize}
\par{\footnotesize\color{mainred}\textbf{代表论文}}
\begin{itemize}[leftmargin=1.2em,itemsep=0pt,topsep=0pt]
\item {\footnotesize Practical techniques for high precision measurements on near-term quantum hardware: a Case Study in Molecular Energy Estimation (npj Quantum Informati}
\item {\footnotesize Improving shadow estimation with locally optimal dual frames (Physical Review A, accepted 2026)}
\end{itemize}
\par{\footnotesize\color{mainred}\textbf{成就}}
\begin{itemize}[leftmargin=1.2em,itemsep=0pt,topsep=0pt]
\item {\footnotesize 在npj Quantum Information发表论文（2025年）}
\item {\footnotesize 在IBM Eagle r3量子硬件上对BODIPY分子进行分子能量估计}
\item {\footnotesize 硕士学位论文：探索在VQE中使用自适应POVM测量（赫尔辛基大学，2022年）}
\end{itemize}
\par{\footnotesize\color{mainred}\textbf{角色}} {\footnotesize 研究员（初级研究员）——为近期量子计算机开发高精度测量技术}
\vspace{2pt}\hrule\vspace{2pt}
\vspace{4pt}
\noindent{\bfseries Matea Leahy}\quad{\color{graytext}\footnotesize Researcher (Junior Researcher / Researcher Intern)}
\par{\footnotesize\color{graytext}MSc in Quantum Science and Technology (ongoing) - Applied and Computational Mathematics (BSc); Quantum Science and Technology (MSc) (Ongoing) | University College Dublin (UCD), Trinity College Dublin (TCD) \textbar\  h-index: 1–2 (very early career) \textbar\  \href{https://www.linkedin.com/in/matea-l-075993100/}{LinkedIn} \textbar\  \href{Not publicly found}{Scholar}}
\par{\footnotesize\color{mainred}\textbf{经历}}
\begin{itemize}[leftmargin=1.2em,itemsep=0pt,topsep=0pt]
\item {\footnotesize Algorithmiq: Researcher Intern / Junior Researcher (2022–present)}
\item {\footnotesize Trinity Quantum Alliance: MSc Thesis Collaborator (2022–present)}
\end{itemize}
\par{\footnotesize\color{mainred}\textbf{研究方向}}
\begin{itemize}[leftmargin=1.2em,itemsep=0pt,topsep=0pt]
\item {\footnotesize 量子错误缓解}
\item {\footnotesize 张量网络方法}
\item {\footnotesize 量子算法}
\end{itemize}
\par{\footnotesize\color{mainred}\textbf{代表论文}}
\begin{itemize}[leftmargin=1.2em,itemsep=0pt,topsep=0pt]
\item {\footnotesize Scalable tensor-network error mitigation for near-term quantum computing (arXiv:2307.11740) -- co-author}
\item {\footnotesize Dynamical simulations of many-body quantum chaos on a quantum computer (2026, EPFL/TCD collaboration)}
\end{itemize}
\par{\footnotesize\color{mainred}\textbf{成就}}
\begin{itemize}[leftmargin=1.2em,itemsep=0pt,topsep=0pt]
\item {\footnotesize 首届Microsoft Ireland Scholarship获得者，用于攻读Trinity College Dublin的量子科学与技术MSc学位（2021/22）}
\item {\footnotesize 在WERQSHOP 2024上展示了可扩展的张量网络错误缓解工作}
\item {\footnotesize 入选Trinity College Dublin Impact Report 2021–22}
\end{itemize}
\par{\footnotesize\color{mainred}\textbf{角色}} {\footnotesize 初级研究员/研究实习生——从事量子错误缓解和算法开发工作，与Algorithmiq的Helsinki团队以及Dublin的Trinity Quantum Alliance合作}
\vspace{2pt}\hrule\vspace{2pt}
\vspace{4pt}
\noindent{\bfseries Joonas Malmi}\quad{\color{graytext}\footnotesize Researcher (Junior Researcher)}
\par{\footnotesize\color{graytext}MSc (PhD student) - Theoretical Physics / Quantum Information (MSc 2020; PhD ongoing) | University of Turku (BSc, MSc), University of Helsinki (PhD, ongoing) \textbar\  h-index: 3–4 \textbar\  \href{https://www.linkedin.com/in/joonas-malmi-1a5466238/}{LinkedIn} \textbar\  \href{Not publicly found (MathSciNet author ID: 1624464)}{Scholar}}
\par{\footnotesize\color{mainred}\textbf{经历}}
\begin{itemize}[leftmargin=1.2em,itemsep=0pt,topsep=0pt]
\item {\footnotesize Algorithmiq Ltd: Junior Researcher (2022–present)}
\item {\footnotesize University of Helsinki — QTF Centre of Excellence: PhD Researcher (2021–present)}
\item {\footnotesize InstituteQ — Finnish Quantum Institute: Researcher (2021–present)}
\end{itemize}
\par{\footnotesize\color{mainred}\textbf{研究方向}}
\begin{itemize}[leftmargin=1.2em,itemsep=0pt,topsep=0pt]
\item {\footnotesize 量子信息理论}
\item {\footnotesize 测量协议（影子估计、对偶框架）}
\item {\footnotesize 复杂网络上的量子行走}
\item {\footnotesize 可观测量估计与经典影子}
\end{itemize}
\par{\footnotesize\color{mainred}\textbf{代表论文}}
\begin{itemize}[leftmargin=1.2em,itemsep=0pt,topsep=0pt]
\item {\footnotesize Enhanced observable estimation through classical optimization of informationally over-complete measurement data — beyond classical shadows (Phys. Rev.}
\item {\footnotesize Improving shadow estimation with locally-optimal dual frames (arXiv:2511.02555, 2025) — co-author}
\end{itemize}
\par{\footnotesize\color{mainred}\textbf{成就}}
\begin{itemize}[leftmargin=1.2em,itemsep=0pt,topsep=0pt]
\item {\footnotesize Physical Review A 和 Physical Review Research 论文的第一作者}
\item {\footnotesize 开发了针对对偶POVM算子的测量后优化方法，在经典影子导致指数开销的情况下实现了零方差估计}
\item {\footnotesize 关于现实世界复杂网络（互联网）上量子空间搜索的硕士学位论文——以PRResearch论文形式发表}
\end{itemize}
\par{\footnotesize\color{mainred}\textbf{角色}} {\footnotesize 初级研究员——从事量子测量协议、影子估计和量子信息理论的研究；多篇带有Algorithmiq单位的论文的通讯作者}
\vspace{2pt}\hrule\vspace{2pt}
\vspace{4pt}
\noindent{\bfseries Aaron Fitzpatrick}\quad{\color{graytext}\footnotesize Researcher (Junior Researcher)}
\par{\footnotesize\color{graytext}PhD-level researcher (formal degree not publicly confirmed) - Quantum Chemistry / Computational Chemistry (Not publicly available) | Trinity Quantum Alliance, Trinity College Dublin \textbar\  h-index: 3–5 \textbar\  \href{https://www.linkedin.com/in/aaron-fitzpatrick-8004931bb/}{LinkedIn} \textbar\  \href{Profile exists at Semantic Scholar (ID: 2065154026) but URL could not be fetched}{Scholar}}
\par{\footnotesize\color{mainred}\textbf{经历}}
\begin{itemize}[leftmargin=1.2em,itemsep=0pt,topsep=0pt]
\item {\footnotesize Algorithmiq Ltd: Junior Researcher (2022–present)}
\item {\footnotesize Trinity Quantum Alliance, TCD: Researcher (Not specified)}
\end{itemize}
\par{\footnotesize\color{mainred}\textbf{研究方向}}
\begin{itemize}[leftmargin=1.2em,itemsep=0pt,topsep=0pt]
\item {\footnotesize 变分量子本征求解器（VQE）方法}
\item {\footnotesize 轨道优化技术}
\item {\footnotesize 用于量子-经典方法的可极化嵌入（PE）}
\item {\footnotesize 近期量子器件的动态电子关联}
\item {\footnotesize 从量子测量重构密度矩阵}
\end{itemize}
\par{\footnotesize\color{mainred}\textbf{代表论文}}
\begin{itemize}[leftmargin=1.2em,itemsep=0pt,topsep=0pt]
\item {\footnotesize Self-Consistent Field Approach for the Variational Quantum Eigensolver: Orbital Optimization Goes Adaptive (J. Phys. Chem. A, 2023) — first author}
\item {\footnotesize The variational quantum eigensolver self-consistent field method within a polarizable embedded framework (J. Chem. Phys. 160, 124114, 2024) — co-autho}
\end{itemize}
\par{\footnotesize\color{mainred}\textbf{成就}}
\begin{itemize}[leftmargin=1.2em,itemsep=0pt,topsep=0pt]
\item {\footnotesize 开发了用于轨道优化的ADAPT-VQE-SCF方法}
\item {\footnotesize 为结合可极化嵌入与量子计算的PE-VQE-SCF方法做出了贡献}
\item {\footnotesize 邮箱: aaron.fitzpatrick@algorithmiq.fi; ORCID: 0009-0007-9921-3522}
\end{itemize}
\par{\footnotesize\color{mainred}\textbf{角色}} {\footnotesize 初级研究员——从事量子化学的混合量子-经典算法研究，特别是具有轨道优化和可极化嵌入的VQE方法}
\vspace{2pt}\hrule\vspace{2pt}
\vspace{4pt}
\noindent{\bfseries Aaron Miller}\quad{\color{graytext}\footnotesize Researcher (Junior Researcher)}
\par{\footnotesize\color{graytext}PhD student-level - Physics / Quantum Information (Not publicly available) | School of Physics, Trinity College Dublin, Trinity Quantum Alliance, Dublin \textbar\  h-index: 4–5 \textbar\  \href{https://www.linkedin.com/in/aaron-miller-07a990197/}{LinkedIn} \textbar\  \href{Not publicly found}{Scholar}}
\par{\footnotesize\color{mainred}\textbf{经历}}
\begin{itemize}[leftmargin=1.2em,itemsep=0pt,topsep=0pt]
\item {\footnotesize Algorithmiq Ltd: Junior Researcher (2022–present)}
\item {\footnotesize School of Physics, Trinity College Dublin: PhD Researcher (Not specified)}
\item {\footnotesize Trinity Quantum Alliance: Researcher (Not specified)}
\end{itemize}
\par{\footnotesize\color{mainred}\textbf{研究方向}}
\begin{itemize}[leftmargin=1.2em,itemsep=0pt,topsep=0pt]
\item {\footnotesize 费米子到量子比特映射（Bonsai、Treespilation算法）}
\item {\footnotesize 量子线路编译}
\item {\footnotesize 变分量子本征求解器（VQE）}
\item {\footnotesize 核-电子轨道（NEO）方法}
\item {\footnotesize 量子处理器上的分子模拟}
\end{itemize}
\par{\footnotesize\color{mainred}\textbf{代表论文}}
\begin{itemize}[leftmargin=1.2em,itemsep=0pt,topsep=0pt]
\item {\footnotesize Bonsai Algorithm: Grow Your Own Fermion-to-Qubit Mappings (PRX Quantum 4, 030314, 2023) — first author}
\item {\footnotesize Treespilation: Architecture- and State-Optimised Fermion-to-Qubit Mappings (arXiv:2403.03992, 2024) — first author}
\end{itemize}
\par{\footnotesize\color{mainred}\textbf{成就}}
\begin{itemize}[leftmargin=1.2em,itemsep=0pt,topsep=0pt]
\item {\footnotesize 开发了用于自动生成优化费米子到量子比特映射的Bonsai算法}
\item {\footnotesize 开发了Treespilation，使ADAPT-VQE的CNOT门计数减少高达74\%}
\item {\footnotesize 与AstraZeneca、IBM Quantum、University of Cambridge合作进行质子转移动力学模拟}
\end{itemize}
\par{\footnotesize\color{mainred}\textbf{角色}} {\footnotesize 初级研究员 — 主导费米子到量子比特映射（Bonsai, Treespilation）及分子模拟量子线路编译工作；同时隶属于Trinity College Dublin}
\vspace{2pt}\hrule\vspace{2pt}
\vspace{4pt}
\noindent{\bfseries Leander Thiessen}\quad{\color{graytext}\footnotesize Researcher (Junior Researcher)}
\par{\footnotesize\color{graytext}Not publicly available - Quantum Chemistry (Not publicly available) | Affiliated with Algorithmiq Ltd, Helsinki, Finland \textbar\  h-index: 1 \textbar\  \href{https://www.linkedin.com/in/leander-thiessen-9285a0113/}{LinkedIn} \textbar\  \href{Not publicly found (INSPIRE author ID: L.Thiessen.1 at https://inspirehep.net/authors/2924717)}{Scholar}}
\par{\footnotesize\color{mainred}\textbf{经历}}
\begin{itemize}[leftmargin=1.2em,itemsep=0pt,topsep=0pt]
\item {\footnotesize Algorithmiq Ltd: Junior Researcher (2024–present)}
\end{itemize}
\par{\footnotesize\color{mainred}\textbf{研究方向}}
\begin{itemize}[leftmargin=1.2em,itemsep=0pt,topsep=0pt]
\item {\footnotesize 激发态量子化学}
\item {\footnotesize ADAPT-VQE方法}
\item {\footnotesize 用于光动力疗法（PDT）的BODIPY光敏剂}
\end{itemize}
\par{\footnotesize\color{mainred}\textbf{代表论文}}
\begin{itemize}[leftmargin=1.2em,itemsep=0pt,topsep=0pt]
\item {\footnotesize Toward Accurate Calculation of Excitation Energies on Quantum Computers with ΔADAPT-VQE: A Case Study of BODIPY Derivatives (J. Phys. Chem. Lett. 15, }
\end{itemize}
\par{\footnotesize\color{mainred}\textbf{成就}}
\begin{itemize}[leftmargin=1.2em,itemsep=0pt,topsep=0pt]
\item {\footnotesize 合著发表了关于BODIPY光敏剂的量子计算方法的论文，激发能的平均绝对误差低至0.054 eV——性能优于经典的TDDFT和EOM-CCSD方法}
\end{itemize}
\par{\footnotesize\color{mainred}\textbf{角色}} {\footnotesize Algorithmiq公司青年研究员——利用VQE方法进行激发态分子模拟的量子化学算法研究}
\vspace{2pt}\hrule\vspace{2pt}
\vspace{4pt}
\noindent{\bfseries Hetta Vappula}\quad{\color{graytext}\footnotesize Junior Researcher}
\par{\footnotesize\color{graytext}Not publicly confirmed (likely BSc/MSc-level) - Quantum Information / Physics (Not publicly available) | University of Helsinki (affiliation inferred from collaborations), Algorithmiq Ltd, Helsinki \textbar\  h-index: 1–2 \textbar\  \href{https://www.linkedin.com/in/hetta-vappula-628055149/}{LinkedIn} \textbar\  \href{Not publicly found (Semantic Scholar author: https://www.semanticscholar.org/author/Hetta-Vappula/2319602989)}{Scholar}}
\par{\footnotesize\color{mainred}\textbf{经历}}
\begin{itemize}[leftmargin=1.2em,itemsep=0pt,topsep=0pt]
\item {\footnotesize Algorithmiq Ltd: Junior Researcher (2024–present)}
\end{itemize}
\par{\footnotesize\color{mainred}\textbf{研究方向}}
\begin{itemize}[leftmargin=1.2em,itemsep=0pt,topsep=0pt]
\item {\footnotesize 近期量子硬件上的高精度测量}
\item {\footnotesize 局部最优对偶框架的阴影估计}
\item {\footnotesize 读出误差缓解}
\item {\footnotesize IBM硬件上的量子化学应用}
\end{itemize}
\par{\footnotesize\color{mainred}\textbf{代表论文}}
\begin{itemize}[leftmargin=1.2em,itemsep=0pt,topsep=0pt]
\item {\footnotesize Practical Techniques for High-Precision Measurements on Near-Term Quantum Hardware: A Case Study in Molecular Energy Estimation (npj Quantum Informati}
\item {\footnotesize Improving shadow estimation with locally-optimal dual frames (arXiv:2511.02555, 2025) — co-author}
\end{itemize}
\par{\footnotesize\color{mainred}\textbf{成就}}
\begin{itemize}[leftmargin=1.2em,itemsep=0pt,topsep=0pt]
\item {\footnotesize 合著《npj Quantum Information》论文，展示了在IBM Eagle r3量子处理器上将测量误差从1–5\%降低至0.16\%}
\item {\footnotesize 对局部最优阴影估计技术做出贡献}
\end{itemize}
\par{\footnotesize\color{mainred}\textbf{角色}} {\footnotesize 初级研究员——致力于近期量子设备的实用误差缓解和高精度测量技术}
\vspace{2pt}\hrule\vspace{2pt}
\vspace{4pt}
\noindent{\bfseries Francesco Grieco}\quad{\color{graytext}\footnotesize Junior Researcher}
\par{\footnotesize\color{graytext}Not publicly available - Not publicly identified (Not publicly available) | Algorithmiq Ltd, Helsinki \textbar\  h-index: 0 \textbar\  \href{https://www.linkedin.com/in/francesco-grieco-58212a346/}{LinkedIn} \textbar\  \href{Not found}{Scholar}}
\par{\footnotesize\color{mainred}\textbf{经历}}
\begin{itemize}[leftmargin=1.2em,itemsep=0pt,topsep=0pt]
\item {\footnotesize Algorithmiq Ltd: Junior Researcher (2024–present)}
\end{itemize}
\par{\footnotesize\color{mainred}\textbf{研究方向}}
\begin{itemize}[leftmargin=1.2em,itemsep=0pt,topsep=0pt]
\item {\footnotesize 未公开确认——未发现任何出版物}
\end{itemize}
\par{\footnotesize\color{mainred}\textbf{代表论文}}
\begin{itemize}[leftmargin=1.2em,itemsep=0pt,topsep=0pt]
\item {\footnotesize No publications found in public databases or search engines}
\end{itemize}
\par{\footnotesize\color{mainred}\textbf{成就}}
\begin{itemize}[leftmargin=1.2em,itemsep=0pt,topsep=0pt]
\item {\footnotesize 在Algorithmiq官方团队页面列为初级研究员}
\end{itemize}
\par{\footnotesize\color{mainred}\textbf{角色}} {\footnotesize 初级研究员——具体研究重点未公开记录}
\vspace{2pt}\hrule\vspace{2pt}
\vspace{4pt}
\noindent{\bfseries Fabio Tarocco}\quad{\color{graytext}\footnotesize Researcher}
\par{\footnotesize\color{graytext}PhD Student (in progress) - Quantum Computation (Ongoing) | University of L'Aquila, University of Verona \textbar\  h-index: 2 (18 total citations) \textbar\  \href{https://www.linkedin.com/in/quantumfabiotarocco/}{LinkedIn} \textbar\  \href{Not publicly indexed; AD Scientific Index profile available at https://adscientificindex.com/scientist/fabio-tarocco/5178248}{Scholar}}
\par{\footnotesize\color{mainred}\textbf{经历}}
\begin{itemize}[leftmargin=1.2em,itemsep=0pt,topsep=0pt]
\item {\footnotesize University of Verona: Master's Student in Engineering and Computer Science (2020-2022)}
\end{itemize}
\par{\footnotesize\color{mainred}\textbf{研究方向}}
\begin{itemize}[leftmargin=1.2em,itemsep=0pt,topsep=0pt]
\item {\footnotesize 变分量子特征求解器（VQE）}
\item {\footnotesize 量子近似优化算法（QAOA）}
\item {\footnotesize 量子退火}
\item {\footnotesize 量子游走与图问题}
\item {\footnotesize 量子化学 - 电子结构与活性空间选择}
\end{itemize}
\par{\footnotesize\color{mainred}\textbf{代表论文}}
\begin{itemize}[leftmargin=1.2em,itemsep=0pt,topsep=0pt]
\item {\footnotesize AEGISS - Atomic orbital and Entropy-based Guided Inference for Space Selection (arXiv:2508.10671, 2025)}
\item {\footnotesize Scalable quantum circuit generation for iterative ground state approximation using Majorana Propagation (arXiv:2603.23444, 2026)}
\end{itemize}
\par{\footnotesize\color{mainred}\textbf{成就}}
\begin{itemize}[leftmargin=1.2em,itemsep=0pt,topsep=0pt]
\item {\footnotesize 作为Algorithmiq团队成员，因量子驱动的癌症药物发现而获得200万美元Wellcome Leap Q4Bio奖（2026年）}
\item {\footnotesize 获得PNRR资助，在University of L'Aquila攻读量子计算博士项目}
\end{itemize}
\par{\footnotesize\color{mainred}\textbf{角色}} {\footnotesize 研究员 - 致力于量子化学与量子计算研究，从事活性空间选择工作流和VQE算法工作}
\vspace{2pt}\hrule\vspace{2pt}
\subsection{工程团队（3人）}
\vspace{4pt}
\noindent{\bfseries Francesca Pietracaprina}\quad{\color{graytext}\footnotesize Performance Engineer}
\par{\footnotesize\color{graytext}PhD - Statistical Physics / Quantum Many-Body Physics (2015) | SISSA (Trieste), Sapienza University of Rome \textbar\  h-index: 11 (21-23 publications, \textasciitilde{}473 total citations) \textbar\  \href{https://www.linkedin.com/in/francesca-pietracaprina/}{LinkedIn} \textbar\  \href{Not directly found; indexed on zbMATH (https://zbmath.org/authors/pietracaprina.francesca), MathSciNet (Author ID 1220640), and arXivLens (https://arxivlens.com/Author/Details/francesca-pietracaprina-2141-25463508)}{Scholar}}
\par{\footnotesize\color{mainred}\textbf{经历}}
\begin{itemize}[leftmargin=1.2em,itemsep=0pt,topsep=0pt]
\item {\footnotesize Sapienza University of Rome: Postdoctoral Researcher (2015-2017)}
\item {\footnotesize University of Toulouse (LPT, CNRS): Postdoctoral Researcher (2017-2019)}
\item {\footnotesize Trinity College Dublin (QUSYS Group): Marie Sklodowska-Curie Fellow (2019-2021)}
\item {\footnotesize Algorithmiq: Software Engineer / Performance Engineer (2022-present)}
\end{itemize}
\par{\footnotesize\color{mainred}\textbf{研究方向}}
\begin{itemize}[leftmargin=1.2em,itemsep=0pt,topsep=0pt]
\item {\footnotesize 无序量子系统中的多体局域化}
\item {\footnotesize Hilbert空间碎裂与多重分形}
\item {\footnotesize 多体系统中的量子纠缠}
\item {\footnotesize Bethe晶格上的Anderson局域化}
\item {\footnotesize 量子热力学}
\end{itemize}
\par{\footnotesize\color{mainred}\textbf{代表论文}}
\begin{itemize}[leftmargin=1.2em,itemsep=0pt,topsep=0pt]
\item {\footnotesize Hilbert-space fragmentation, multifractality, and many-body localization - Annals of Physics (2021)}
\item {\footnotesize Probing many-body localization in a disordered quantum dimer model on the honeycomb lattice - SciPost Phys. 10, 044 (2021)}
\end{itemize}
\par{\footnotesize\color{mainred}\textbf{成就}}
\begin{itemize}[leftmargin=1.2em,itemsep=0pt,topsep=0pt]
\item {\footnotesize 都柏林圣三一学院玛丽·居里学者（欧盟Horizon 2020计划）}
\item {\footnotesize 参与获得欧洲研究委员会概念验证资助的研究项目，将统计物理学工具应用于卫星图像分析以评估荒漠化风险}
\item {\footnotesize 2024年林道诺贝尔奖得主大会参会者}
\end{itemize}
\par{\footnotesize\color{mainred}\textbf{角色}} {\footnotesize 性能工程师（亦列为软件工程师与研究员）——专注于优化量子算法性能及全栈量子计算解决方案，以应用于药物发现与模拟领域}
\vspace{2pt}\hrule\vspace{2pt}
\vspace{4pt}
\noindent{\bfseries Ramon Panades Barrueta}\quad{\color{graytext}\footnotesize Software Engineer}
\par{\footnotesize\color{graytext}PhD in Physics - Quantum Chemistry / Theoretical Physics (2020) | Universite de Lille, University of Havana \textbar\  h-index: 1-2 (5 publications, \textasciitilde{}16 citations) \textbar\  \href{https://www.linkedin.com/in/rpanades/}{LinkedIn} \textbar\  \href{https://scholar.google.com/citations?user=IMWuq6wAAAAJ&hl}{Scholar}}
\par{\footnotesize\color{mainred}\textbf{经历}}
\begin{itemize}[leftmargin=1.2em,itemsep=0pt,topsep=0pt]
\item {\footnotesize TU Dresden (NOMAD Laboratory): Postdoctoral Researcher (2021-2023)}
\item {\footnotesize University of Twente (TREX Project): Postdoctoral Researcher (2020-2021)}
\item {\footnotesize Heidelberg University: Visiting PhD Student (2019)}
\item {\footnotesize University of Lille 1: Teaching Assistant (2018-2019)}
\end{itemize}
\par{\footnotesize\color{mainred}\textbf{研究方向}}
\begin{itemize}[leftmargin=1.2em,itemsep=0pt,topsep=0pt]
\item {\footnotesize 量子计算与张量网络算法（混合量子-经典）}
\item {\footnotesize 在FHI-aims和GreenX中实现的用于核心级光谱学（XPS/XAS）的低标度GW基方法}
\item {\footnotesize 大规模并行量子蒙特卡洛算法（CHAMP代码，TREX CoE）}
\item {\footnotesize 核量子动力学与张量分解（MCTDH方法）}
\item {\footnotesize 势能面生成（SRP-MGPF、SOP-FBR、CP-FBR方法）}
\end{itemize}
\par{\footnotesize\color{mainred}\textbf{代表论文}}
\begin{itemize}[leftmargin=1.2em,itemsep=0pt,topsep=0pt]
\item {\footnotesize Accelerating core-level GW calculations by combining the contour deformation approach with the analytic continuation of W - J. Chem. Theory Comput. (2}
\item {\footnotesize Double excitation energies from quantum Monte Carlo using state-specific energy optimization - J. Chem. Theory Comput. (2022, with Shepard et al.)}
\end{itemize}
\par{\footnotesize\color{mainred}\textbf{成就}}
\begin{itemize}[leftmargin=1.2em,itemsep=0pt,topsep=0pt]
\item {\footnotesize 博士论文：大气分子与模型烟尘颗粒相互作用的量子模拟（Universite de Lille）}
\item {\footnotesize 用于自动化势能面生成的SRPTucker软件包的开发者}
\item {\footnotesize 开源量子化学代码（FHI-aims、GreenX、CHAMP）的贡献者}
\end{itemize}
\par{\footnotesize\color{mainred}\textbf{角色}} {\footnotesize 软件工程师 - 使用张量网络方法实现前沿的高性能混合量子-经典算法；致力于编码最佳实践和稳健的软件设计}
\vspace{2pt}\hrule\vspace{2pt}
\vspace{4pt}
\noindent{\bfseries Cannur Floszmann}\quad{\color{graytext}\footnotesize Junior e-Learning Developer}
\par{\footnotesize\color{graytext}Not publicly available - Not publicly available (Not publicly available) \textbar\  h-index: N/A \textbar\  \href{https://www.linkedin.com/in/cannur-floszmann/}{LinkedIn} \textbar\  \href{Not found}{Scholar}}
\par{\footnotesize\color{mainred}\textbf{角色}} {\footnotesize 初级电子学习开发员——负责开发电子学习内容和教育材料（可能与量子计算教育平台相关）}
\vspace{2pt}\hrule\vspace{2pt}
\subsection{商业与运营（3人）}
\vspace{4pt}
\noindent{\bfseries Elsi-Mari Borrelli}\quad{\color{graytext}\footnotesize Country Manager, Finland}
\par{\footnotesize\color{graytext}PhD in Theoretical Physics - Theoretical Physics / Quantum Physics (Not publicly specified) | University of Helsinki (presumed), Aalto University (postdoc) \textbar\  h-index: \textasciitilde{}15 (estimated based on highly-cited papers in non \textbar\  \href{https://www.linkedin.com/in/elsi-mari-borrelli/}{LinkedIn} \textbar\  \href{https://scholar.google.com/citations?user=J9G1CR8AAAAJ}{Scholar}}
\par{\footnotesize\color{mainred}\textbf{经历}}
\begin{itemize}[leftmargin=1.2em,itemsep=0pt,topsep=0pt]
\item {\footnotesize ABB Corporate Research Switzerland: Senior Scientist (\textasciitilde{}2018-2022)}
\item {\footnotesize Aalto University: Postdoctoral Researcher ()}
\item {\footnotesize University of Turku: Postdoctoral Researcher ()}
\end{itemize}
\par{\footnotesize\color{mainred}\textbf{研究方向}}
\begin{itemize}[leftmargin=1.2em,itemsep=0pt,topsep=0pt]
\item {\footnotesize 量子生物学与量子网络医学}
\item {\footnotesize 光动力疗法（PDT）光敏剂设计}
\item {\footnotesize 量子错误缓解与测量断层成像}
\item {\footnotesize 非马尔可夫开放量子系统}
\item {\footnotesize ADAPT-VQE变分量子算法}
\end{itemize}
\par{\footnotesize\color{mainred}\textbf{代表论文}}
\begin{itemize}[leftmargin=1.2em,itemsep=0pt,topsep=0pt]
\item {\footnotesize Colloquium: Non-Markovian dynamics in open quantum systems - Reviews of Modern Physics (2016) - 800+ citations}
\item {\footnotesize Measure for the degree of non-Markovian behavior - Physical Review Letters (2009) - 1,300+ citations}
\end{itemize}
\par{\footnotesize\color{mainred}\textbf{成就}}
\begin{itemize}[leftmargin=1.2em,itemsep=0pt,topsep=0pt]
\item {\footnotesize 合著了《Reviews of Modern Physics》中关于非马尔可夫动力学的开创性座谈会论文}
\item {\footnotesize 量子网络医学新兴领域的先驱}
\item {\footnotesize Algorithmiq量子生物学和光动力疗法研究项目的主要贡献者}
\end{itemize}
\par{\footnotesize\color{mainred}\textbf{角色}} {\footnotesize 芬兰国家经理（同时在Algorithmiq团队页面上列为企业合作与销售主管；曾任首席量子生物学研究员）}
\vspace{2pt}\hrule\vspace{2pt}
\vspace{4pt}
\noindent{\bfseries David Lin}\quad{\color{graytext}\footnotesize General Counsel}
\par{\footnotesize\color{graytext}Juris Doctor (J.D.) - Electrical Engineering / Law (Not publicly specified) | UCLA - B.S. in Electrical Engineering, Loyola Law School, Los Angeles - J.D. \textbar\  h-index: N/A \textbar\  \href{https://www.linkedin.com/in/dklin/}{LinkedIn} \textbar\  \href{N/A}{Scholar}}
\par{\footnotesize\color{mainred}\textbf{经历}}
\begin{itemize}[leftmargin=1.2em,itemsep=0pt,topsep=0pt]
\item {\footnotesize Winston \& Strawn LLP: BigLaw Litigator (2013-2023)}
\item {\footnotesize davidlin.law, PC: Founder ()}
\item {\footnotesize DTO Law (Los Angeles): Counsel (Current)}
\end{itemize}
\par{\footnotesize\color{mainred}\textbf{成就}}
\begin{itemize}[leftmargin=1.2em,itemsep=0pt,topsep=0pt]
\item {\footnotesize 代表三家主要视频游戏发行商在PTAB前成功无效了一个专利组合}
\item {\footnotesize 在平行的地区法院诉讼中,通过简易判决成功击败了六项被主张的专利}
\item {\footnotesize 在国际贸易委员会(ITC)审理了一桩涉及多项专利的案件,经历了完整的庭审过程}
\end{itemize}
\par{\footnotesize\color{mainred}\textbf{角色}} {\footnotesize 法务总监——负责知识产权战略、技术商业交易和合作协议。同时担任DTO Law的法律顾问。常驻加利福尼亚州帕萨迪纳。}
\vspace{2pt}\hrule\vspace{2pt}
\vspace{4pt}
\noindent{\bfseries Jason Gumarang}\quad{\color{graytext}\footnotesize IT Specialist}
\par{\footnotesize\color{graytext}Not publicly available - Not publicly available (Not publicly available) \textbar\  h-index: N/A \textbar\  \href{https://www.linkedin.com/in/jason-gumarang/}{LinkedIn} \textbar\  \href{N/A}{Scholar}}
\par{\footnotesize\color{mainred}\textbf{经历}}
\begin{itemize}[leftmargin=1.2em,itemsep=0pt,topsep=0pt]
\item {\footnotesize Not publicly available: Not publicly available (Not publicly available)}
\end{itemize}
\par{\footnotesize\color{mainred}\textbf{成就}}
\begin{itemize}[leftmargin=1.2em,itemsep=0pt,topsep=0pt]
\item {\footnotesize 在 Algorithmiq 官方团队页面列为 IT 专员}
\end{itemize}
\par{\footnotesize\color{mainred}\textbf{角色}} {\footnotesize IT 专员 - 负责 IT 基础设施和系统支持}
\vspace{2pt}\hrule\vspace{2pt}
\section{研究方向分布}
\begin{center}\begin{tikzpicture}[scale=0.85]
\fill[mainred] (0,0.00) rectangle (9.50,0.55);
\draw[gray!50] (0,0.00) rectangle (9.50,0.55);
\node[anchor=east] at (-0.3,0.28) {\footnotesize 量子算法};
\node[anchor=west] at (9.70,0.28) {\footnotesize\color{graytext} 14};
\fill[mainred] (0,-0.85) rectangle (8.82,-0.30);
\draw[gray!50] (0,-0.85) rectangle (8.82,-0.30);
\node[anchor=east] at (-0.3,-0.57) {\footnotesize 量子化学};
\node[anchor=west] at (9.02,-0.57) {\footnotesize\color{graytext} 13};
\fill[mainred] (0,-1.70) rectangle (8.14,-1.15);
\draw[gray!50] (0,-1.70) rectangle (8.14,-1.15);
\node[anchor=east] at (-0.3,-1.42) {\footnotesize 量子模拟};
\node[anchor=west] at (8.34,-1.42) {\footnotesize\color{graytext} 12};
\fill[mainred] (0,-2.55) rectangle (6.11,-2.00);
\draw[gray!50] (0,-2.55) rectangle (6.11,-2.00);
\node[anchor=east] at (-0.3,-2.27) {\footnotesize 量子信息};
\node[anchor=west] at (6.31,-2.27) {\footnotesize\color{graytext} 9};
\fill[mainred] (0,-3.40) rectangle (3.39,-2.85);
\draw[gray!50] (0,-3.40) rectangle (3.39,-2.85);
\node[anchor=east] at (-0.3,-3.12) {\footnotesize 误差缓解};
\node[anchor=west] at (3.59,-3.12) {\footnotesize\color{graytext} 5};
\fill[mainred] (0,-4.25) rectangle (2.71,-3.70);
\draw[gray!50] (0,-4.25) rectangle (2.71,-3.70);
\node[anchor=east] at (-0.3,-3.98) {\footnotesize 张量网络};
\node[anchor=west] at (2.91,-3.98) {\footnotesize\color{graytext} 4};
\fill[mainred] (0,-5.10) rectangle (2.04,-4.55);
\draw[gray!50] (0,-5.10) rectangle (2.04,-4.55);
\node[anchor=east] at (-0.3,-4.82) {\footnotesize 机器学习};
\node[anchor=west] at (2.24,-4.82) {\footnotesize\color{graytext} 3};
\fill[mainred] (0,-5.95) rectangle (1.36,-5.40);
\draw[gray!50] (0,-5.95) rectangle (1.36,-5.40);
\node[anchor=east] at (-0.3,-5.67) {\footnotesize 量子生物学};
\node[anchor=west] at (1.56,-5.67) {\footnotesize\color{graytext} 2};
\end{tikzpicture}\end{center}
\section{期刊分布}
\begin{center}\begin{tikzpicture}[scale=0.85]
\fill[mainred] (0,0.00) rectangle (9.50,0.55);
\draw[gray!50] (0,0.00) rectangle (9.50,0.55);
\node[anchor=east] at (-0.3,0.28) {\footnotesize Phys. Rev.};
\node[anchor=west] at (9.70,0.28) {\footnotesize\color{graytext} 16};
\fill[mainred] (0,-0.85) rectangle (3.56,-0.30);
\draw[gray!50] (0,-0.85) rectangle (3.56,-0.30);
\node[anchor=east] at (-0.3,-0.57) {\footnotesize ACS};
\node[anchor=west] at (3.76,-0.57) {\footnotesize\color{graytext} 6};
\fill[mainred] (0,-1.70) rectangle (2.97,-1.15);
\draw[gray!50] (0,-1.70) rectangle (2.97,-1.15);
\node[anchor=east] at (-0.3,-1.42) {\footnotesize Quantum};
\node[anchor=west] at (3.17,-1.42) {\footnotesize\color{graytext} 5};
\fill[mainred] (0,-2.55) rectangle (2.38,-2.00);
\draw[gray!50] (0,-2.55) rectangle (2.38,-2.00);
\node[anchor=east] at (-0.3,-2.27) {\footnotesize Nature};
\node[anchor=west] at (2.58,-2.27) {\footnotesize\color{graytext} 4};
\fill[mainred] (0,-3.40) rectangle (1.78,-2.85);
\draw[gray!50] (0,-3.40) rectangle (1.78,-2.85);
\node[anchor=east] at (-0.3,-3.12) {\footnotesize J. Chem. Phys.};
\node[anchor=west] at (1.98,-3.12) {\footnotesize\color{graytext} 3};
\fill[mainred] (0,-4.25) rectangle (1.19,-3.70);
\draw[gray!50] (0,-4.25) rectangle (1.19,-3.70);
\node[anchor=east] at (-0.3,-3.98) {\footnotesize RSC};
\node[anchor=west] at (1.39,-3.98) {\footnotesize\color{graytext} 2};
\fill[mainred] (0,-5.10) rectangle (0.59,-4.55);
\draw[gray!50] (0,-5.10) rectangle (0.59,-4.55);
\node[anchor=east] at (-0.3,-4.82) {\footnotesize MDPI};
\node[anchor=west] at (0.79,-4.82) {\footnotesize\color{graytext} 1};
\fill[mainred] (0,-5.95) rectangle (0.59,-5.40);
\draw[gray!50] (0,-5.95) rectangle (0.59,-5.40);
\node[anchor=east] at (-0.3,-5.67) {\footnotesize Springer};
\node[anchor=west] at (0.79,-5.67) {\footnotesize\color{graytext} 1};
\fill[mainred] (0,-6.80) rectangle (0.59,-6.25);
\draw[gray!50] (0,-6.80) rectangle (0.59,-6.25);
\node[anchor=east] at (-0.3,-6.52) {\footnotesize World Sci.};
\node[anchor=west] at (0.79,-6.52) {\footnotesize\color{graytext} 1};
\fill[mainred] (0,-7.65) rectangle (0.59,-7.10);
\draw[gray!50] (0,-7.65) rectangle (0.59,-7.10);
\node[anchor=east] at (-0.3,-7.37) {\footnotesize Elsevier};
\node[anchor=west] at (0.79,-7.37) {\footnotesize\color{graytext} 1};
\end{tikzpicture}\end{center}